
\documentclass[a4paper, 20pt,reqno]{article}
\usepackage[left=2cm,right=2cm,top=2cm,bottom=2cm]{geometry}
\usepackage[T2A]{fontenc}
% \DeclareMathSizes{10}{10}{10}{10}

\usepackage[russian]{babel}
\usepackage{soul}
\usepackage{gensymb}
\usepackage{amsfonts,amsmath,amssymb}
\usepackage{mathrsfs}
\usepackage{graphicx}
\usepackage[normalem]{ulem}
\usepackage[document]{ragged2e}
\usepackage{stmaryrd}
\usepackage{wrapfig}
\usepackage{fancyhdr}
\usepackage{floatflt}
\usepackage{python}
\usepackage{float}
\usepackage{amssymb}
\usepackage[most]{tcolorbox}
\usepackage{indentfirst}
\usepackage{setspace}
\usepackage{scrextend}
\usepackage{listings}
\usepackage{makecell,tabularx}
\usepackage{hyperref}
\usepackage{xcolor}

\newcommand{\mycopyright}{pluttan}
\newcommand{\docopyright}{$\mathfrak{Copyright}\ \mathfrak{\mycopyright} \logo$}
\newcommand{\rub}{{\rm{Р}\kern-.635em\rule[.5ex]{.52em}{.04em}\kern.11em}}

\definecolor{linkcolor}{HTML}{000000} 
\definecolor{urlcolor}{HTML}{0000FF} 

\hypersetup{pdfstartview=FitH,  linkcolor=linkcolor,urlcolor=urlcolor, colorlinks=true}

\definecolor{grey}{RGB}{40, 40, 40}

\renewcommand{\href}[1]{\url{#1}}

\definecolor{codegreen}{rgb}{0,0.6,0}
\definecolor{codegray}{rgb}{0.5,0.5,0.5}
\definecolor{codepurple}{rgb}{0.7,0,0.82}
\definecolor{mermaid}{rgb}{0.5,0.3,1}
\definecolor{backcolour}{rgb}{0.95,0.95,0.92}

\lstdefinestyle{mystyle}{
    backgroundcolor=\color{backcolour},   
    commentstyle=\color{codegreen},
    keywordstyle=\color{mermaid},
    numberstyle=\tiny\color{codegray},
    stringstyle=\color{codepurple},
    basicstyle=\ttfamily\footnotesize,
    breakatwhitespace=false,         
    breaklines=true,                 
    captionpos=b,                    
    keepspaces=true,                 
    numbers=left,                    
    numbersep=5pt,                  
    showspaces=false,                
    showstringspaces=false,
    showtabs=false,                  
    tabsize=2,
  extendedchars=true , % включаем не латиницу 
  escapechar=|, % |«выпадаем» в LATEX|
  frame=tb , % рамка сверху и снизу 
  commentstyle=\itshape , % шрифт для комментариев 
  stringstyle=\bfseries
}

\lstset{style=mystyle}

\lstdefinestyle{CommentStyle}{
    language=XML,
    %numbers=left, numberstyle=\tiny, stepnumber=1, numbersep=5pt,
    commentstyle=\color{red},
	basicstyle=\footnotesize\ttfamily,
	language={[ANSI]C++},
	keywordstyle=\bfseries,
	showstringspaces=false,
	morekeywords={include, printf},
	commentstyle={},
	escapeinside=§§,
	escapebegin=\begin{russian}\commentfont,
	escapeend=\end{russian},
    keywordstyle=\color{blue}\bfseries,
    morekeywords={align,begin},
    extendedchars=\true,
    tabsize=2
}
\lstdefinestyle{myLatexStyle}{
    language=c++,
    %backgroundcolor=\color{grey},
    numbers=left, numberstyle=\tiny, stepnumber=1, numbersep=5pt,
    commentstyle=\color{red},
    keywordstyle=\color{blue}\bfseries,
    morekeywords={align,begin},
    extendedchars=\true,
    tabsize=2
}

\lstdefinestyle{pmyLatexStyle}{
    language=java,
    %backgroundcolor=\color{grey},
    numbers=left, numberstyle=\tiny, stepnumber=1, numbersep=5pt,
    commentstyle=\color{red},
    keywordstyle=\color{blue}\bfseries,
    morekeywords={align,begin},
    extendedchars=\true,
    tabsize=2
}

\definecolor{block-gray}{gray}{0.90} % уровень прозрачности (1 - максимум)
\definecolor{yellow}{HTML}{F0FFFF}
\newtcolorbox{myquote}{colback=block-gray,grow to right by=-10mm,grow to left by=-10mm,boxrule=0pt,boxsep=0pt,breakable} % настройки области с изменённым фоном
\newtcolorbox{myquote2}{colback=yellow,grow to right by=-10mm,grow to left by=-10mm,boxrule=0pt,boxsep=0pt,breakable} % настройки области с изменённым фоном

\setlength{\parindent}{12,5mm}

\newcommand{\logo}{\vcenter{\hbox{\includegraphics[width=.8em]{/Users/pluttan/Documents/bw2.png}}}}
\onehalfspacing

\pagestyle{fancy}
\renewcommand{\sectionmark}[1]{\markright{#1}}
\fancyhf{} 
\fancyhead[R]{\bfseries\thepage}
\fancyhead[LO]{\docopyright}

\newcommand{\image}[2]{
	\begin{figure}[H]
		\center{\includegraphics[height=#2pt]{../img/#1} }
    \end{figure}
}

\newcommand{\dotitle}[2]{

\thispagestyle{empty}
\sloppy{
  \scriptsize{
    \line(6,0){0}

    \centering \docopyright

    \centering Привет! Меня зовут Андрей, я создаю свою ботву, этот файл малая ее часть. 
    
    \centering Пользоваться и распространять файлы конечно же можно. Если вы нашли ошибку в файле, можете 
    
    \centering исправить ее в исходном коде и подать на слияние или просто написать в issue. 

    \centering {\bf{Так же вы можете купить распечатанную версию данного файла в виде книжки.}
    
    \centering По всем вопросам писать в ВК.}
    
    \centering Приятного бота)

    \line(6,0){300}

    \centering GitHub: \href{https://github.com/pluttan}

    \centering VK: \href{https://vk.com/pluttan}

    \line(6,0){0}
  }}

\begin{figure}[H]
		\center{\includegraphics[height=100pt]{/Users/pluttan/Documents/forMyDocs/qr2.png} }
\end{figure}

\topskip=-200pt
\vspace*{130pt}
\begin{Huge}
  \textbf{
    \begin{center}
        #1
    \end{center}
  }
  {\begin{center} 
        #2
    \end{center}}
\end{Huge}
\vspace*{300pt}
  \begin{flushright}
    Над файлом работали: \\
    \mycopyright
  \end{flushright}
\newpage
\pagestyle{fancy}
\renewcommand{\sectionmark}[1]{\markright{#1}}
\fancyhf{} 
\fancyhead[R]{\bfseries\thepage}
\fancyhead[LO]{\docopyright}
}
\renewcommand{\sectionmark}[1]{\markright{#1}}
\renewcommand{\sectionmark}[1]{\markright{#1}}
% \makeatletter
%      \renewcommand*\l@section{\@dottedtocline{1}{0em}{1.8em}}
%      \renewcommand*\l@subsection{\@dottedtocline{2}{1.5em}{2.0em}}
%      \renewcommand*\l@subsubsection{\@dottedtocline{3}{4.3em}{3.0em}}
% \makeatother
\newcommand{\toc}{
\newpage
\renewcommand{\contentsname}{Оглавление}
\large{\tableofcontents}
\newpage
\large
}


% It's XeLaTeX, if you can't compile comment this line 
\usepackage{fontspec}\setmainfont{Times New Roman} 

\renewcommand{\copyright}{fiixii, pluttan}
\begin{document}    
    
\dotitle{Подготовка к экзамену}{Математический анализ}
\toc 
\section*{Определения и понятия}\addcontentsline{toc}{section}{Определения и понятия}
\begin{enumerate}

\item $\mathbb{N}$ - \textbf{множество натуральных чисел}, состоит из чисел, возникающих при счёте.
\item $\mathbb{Z}$ - \textbf{множество целых чисел}, состоит из натуральных чисел, нуля и чисел, противоположных натуральным.
\item $\mathbb{Q}$ - \textbf{множество рациональных чисел}, состоит из чисел, представимых в виде ${\dfrac{z}{n}}, z \in \mathbb{Z}, \ n \in \mathbb{N}$.
\item $\mathbb{I}$ - \textbf{множество иррациональных чисел}, состоит из чисел, которые не представимы в виде ${\dfrac{z}{n}}, \ z \in \mathbb{Z}, \ n \in \mathbb{N}$, такие как $e, \pi, \sqrt{3}$ и т.д..
\item $\mathbb{R}$ - \textbf{множество действительных чисел}, состоит из рациональных и иррациональных чисел.
\item $\overline{\mathbb{R}}$ - \textbf{расширенное множество действительных чисел}, состоит из действительных чисел с добавлением $\{+\infty\}$ и $\{-\infty\}$.
\item \textbf{Окрестностью} $U(x)$ \textbf{точки $x$} называют любой интервал, содержащий эту точку.
\item \textbf{Проколотой окрестностью $
\overset{\circ}U(x)$ точки $x$} называют окрестность этой точки $U(x)$, за исключением самой точки $x$.
\item \textbf{${\varepsilon}$-окрестностью точки $x_0$} (при положительном ${\varepsilon}$) называют интервал $(x_0-{\varepsilon}, \ x_0+{\varepsilon}$).

$$U_{\varepsilon}(x_0) = \{x \in \mathbb{R}: \ x_0 - {\varepsilon} < x < x_0 + {\varepsilon}\}
    $$
\item \textbf{Правой (правосторонней) $\delta$-окрестностью точки} $x_0$ называют полуинтервал \ $[x_0,\ x_0 + \delta), \ \delta > 0$.

$$U^+_{\delta}(x_0) = \{x \in \mathbb{R}: \ x_0 \leqslant x < x_0 + {\delta}\}, \ \delta > 0
    $$
\item \textbf{Левой (левосторонней) $\delta$-окрестностью точки} $x_0$ называют полуинтервал \ $(x_0 - \delta,\ x_0], \ \delta > 0$.

$$U^-_{\delta}(x_0) = \{x \in \mathbb{R}: \ x_0 -\delta < x \leqslant x_0\},\ \delta > 0
    $$
\item \textbf{Окрестностью точки} $+\infty$ называют интервал $(a, +\infty), \ a > 0$.

$$U(+\infty) = \{x \in \mathbb{R}: \ x > a\}, \ a > 0
    $$
\item \textbf{Окрестностью точки} $-\infty$ называют интервал $(-\infty, -a), \ a > 0$.

$$U(-\infty) = \{x \in \mathbb{R}: \ x < -a\}, \ a > 0
    $$
\item \textbf{Окрестностью} $\infty$ (бесконечности без знака) называют объединение двух интервалов $(-\infty, -a) \cup (a, +\infty),\ a > 0$.

$$U(\infty) = \{x \in \mathbb{R}: \ | x | > a\}, \ a > 0
    $$
\item \textbf{Последовательностью $\{X_n\}$} называется числовая функция натурального аргумента. Если натуральному числу $n$ при этом поставлено в соответствие число $x_n$, то это число называется $n$-м элементом последовательности; $n$ называют номером элемента $x_n$.
\item Последовательность чисел $\{X_n\}$ называется \textbf{неубывающей}, если $x_{n+1}\geqslant x_n, \ \forall\ n \in \mathbb{N}$.
\item Последовательность чисел $\{X_n\}$ называется \textbf{возрастающей}, если $x_{n+1} > x_n, \ \forall\ n \in \mathbb{N}$.
\item Последовательность чисел $\{X_n\}$ называется \textbf{невозрастающей}, если $x_{n+1} \leqslant x_n, \ \forall\ n \in \mathbb{N}$.
\item Последовательность чисел $\{X_n\}$ называется \textbf{убывающей}, если $x_{n+1} < x_n, \ \forall\ n \in \mathbb{N}$.
\item Неубывающие, невозрастающие, убывающие и возрастающие последовательности называют \textbf{монотонными}.
\item Последовательность называется \textbf{постоянной}, если $ \forall n \in \mathbb{N}: \ x_n = c, \ c \in \mathbb{R}$.
\item Последовательность $\{X_n\}$ называется \textbf{ограниченной сверху}, если $\exists M \in \mathbb{R}$, такое, что $\forall n \in \mathbb{N}: \ x_n \leqslant M$.
\item Последовательность $\{X_n\}$ называется \textbf{ограниченной снизу}, если $\exists M \in \mathbb{R}$, такое, что $\forall n \in \mathbb{N}: \ x_n \geqslant M$.
\item Последовательность, ограниченная и сверху и снизу, называют \textbf{ограниченной}: $\exists M > 0, \ M \in \mathbb{R}$, такое, что $\forall n \in \mathbb{N}: \ \mid  x_n\mid   \leqslant M$.
\item Число $a$ называется \textbf{пределом числовой последовательности} $\{X_n\}$, если для любого, сколь угодно малого положительного ${\varepsilon}$ существует такой номер $N$, зависящий от ${\varepsilon}$, что для всех $n > N$ выполняется неравенство $\mid  a - x_n\mid   < {\varepsilon}$.

$$\lim\limits_{n \to \infty}x_n = a \iff \forall {\varepsilon}>0 \ \ \exists N = N({\varepsilon}) \in \mathbb{N}: \ \forall n > N \Rightarrow | x_n - a|  < {\varepsilon}
    $$
\item Числовая последовательность называется \textbf{сходящейся}, если существует предел этой последовательности, и он конечен.
\item Последовательность $\{X_n\}$ называется \textbf{фундаментальной}, если для любого ${\varepsilon} > 0$ существует номер $N = N({\varepsilon})$ такой, что при любых $m \geqslant N$ и $n \geqslant N$ выполняется неравенство $\mid  x_m - x_n\mid   < {\varepsilon}$.
\item Число $a$ называется \textbf{пределом функции} $f(x)$ при $x \rightarrow x_0$, если для любого ${\varepsilon} > 0$ существует положительное число $\delta = \delta ({\varepsilon})$ такое, что для любого $x \in \overset{\circ}U_\delta(x_0)$ выполняется неравенство $\mid  f(x) - a\mid   < {\varepsilon}$ (определение по Коши).

$$\lim\limits_{x \to x_0} f(x) = a \iff \forall{\varepsilon} > 0 \ \ \exists \delta = \delta({\varepsilon}) > 0 : \ \forall x \in \overset{\circ}U_\delta(x_0) \Rightarrow | f(x) - a|  < {\varepsilon}
    $$
\item Число $a$ называется \textbf{пределом функции} $f(x)$ при $x \rightarrow x_0$, если для любой последовательности $\{X_n\}$ точек из $\overset{\circ}U(x_0)$, для которой $\lim\limits_{n \to \infty}x_n = x_0$, выполняется равенство $\lim\limits_{n \to \infty}\{f(x_n)\} = a$ (определение по Гейне).

$$\lim\limits_{x \to x_0}f(x) = a \iff \{\forall x_n \in \overset{\circ}U(x_0), \ n \in \mathbb{N}\} \ \cap \ \lim\limits_{n \to \infty}{x_n} = x_0 : \ 
    \lim\limits_{n \to \infty}\{f(x_n)\} = a
    $$
\item Число $a$ называется \textbf{правым (правосторонним) пределом функции} f(x) при $x \rightarrow x_0 +$, если для любого ${\varepsilon} > 0$ существует $\delta = \delta({\varepsilon}) > 0$ такое, что при любом $x \in {\overset{\circ\hspace{22pt}}{U^+_\delta(x_0)}} \ \ (\text{т.е.} \ x_0 < x < x_0 + \delta)$, выполняется неравенство $\mid  f(x) - a\mid   < {\varepsilon}$.

$$\lim\limits_{x \to x_0 +} f(x) = a \iff \forall{\varepsilon} > 0 \ \ \exists \delta = \delta({\varepsilon}) > 0 : \ \forall x \in {\overset{\circ\hspace{22pt}}{U^+_\delta(x_0)}}  \Rightarrow | f(x) - a|  < {\varepsilon}
    $$
\item Число $a$ называется \textbf{левым (левосторонним) пределом функции} f(x) при $x \rightarrow x_0 -$, если для любого ${\varepsilon} > 0$ существует $\delta = \delta({\varepsilon}) > 0$ такое, что при любом $x \in {\overset{\circ\hspace{22pt}}{U^-_\delta(x_0)}}, \ \ (\text{т.е.} \ x_0 - \delta < x < x_0)$ выполняется неравенство $\mid  f(x) - a\mid   < {\varepsilon}$.

$$\lim\limits_{x \to x_0 -} f(x) = a \iff \forall{\varepsilon} > 0 \ \ \exists \delta = \delta({\varepsilon}) > 0 : \ \forall x \in {\overset{\circ\hspace{22pt}}{U^-_\delta(x_0)}}  \Rightarrow | f(x) - a|  < {\varepsilon}
    $$
\item Функцию $f(x)$ называют \textbf{ограниченной на множестве $D$}, если существует такое число $M > 0$, что для любых $x \in D$ выполняется неравенство $\mid  f(x)\mid   \leqslant M$.
\item Функцию $f(x)$ называют \textbf{ограниченной} (на области определения $D_f$), если существует такое число $M > 0$, что для любых $x \in D_f$ выполняется неравенство $\mid  f(x)\mid   \leqslant M$.
\item Функцию $f(x)$ называют \textbf{локально ограниченной в окрестности точки $a$}, если существует такое число $M > 0$ и такая окрестность $\overset{\circ}U_\delta(a)$, что для любых $x \in \overset{\circ}U_\delta(a)$ выполняется неравенство $\mid  f(x)\mid   \leqslant M$.
\item Функцию $f(x)$ называют \textbf{бесконечно малой} при $x \rightarrow x_0, \ x_0 \in \overline{\mathbb{R}}$, если $\lim\limits_{x \to x_0} f(x) = 0$.
\item Функцию $f(x)$ называют \textbf{бесконечно большой} при $x \rightarrow x_0, \ x_0 \in \overline{\mathbb{R}}$, если $\lim\limits_{x \to x_0} f(x) = \infty$.
\item $\lim\limits_{x \to 0} {\dfrac{sin(x)}{x}} = 1$ - \textbf{первый замечательный предел}.
\item $\lim\limits_{x \to \infty} \bigg(1 + {\dfrac{1}{x}}\bigg)^{ x}$ - \textbf{второй замечательный предел}.
\item Функции $\alpha(x)$ и $\beta(x)$ называют \textbf{сравнимыми} бесконечно малыми при $x \rightarrow x_0$, если существует хотя бы один из пределов $\lim\limits_{x \to x_0} { \dfrac{\alpha(x)}{\beta(x)}}$ или $\lim\limits_{x \to x_0} { \dfrac{\beta(x)}{\alpha(x)}}$.
\item Функции $\alpha(x)$ и $\beta(x)$ называют \textbf{несравнимыми} бесконечно малыми при $x \rightarrow x_0$, если не существует ни конечного, ни бесконечного предела их отношения при $x \rightarrow x_0$.
\item Функции $\alpha(x)$ и $\beta(x)$ называют \textbf{бесконечно малыми одного порядка} при $x \rightarrow x_0$ и записывают $\alpha(x) = O(\beta(x))$, если существует отличный от нуля конечный предел отношения $\alpha(x)/\beta(x)$, при $x \rightarrow x_0$.

$$\alpha(x) = O(\beta(x)) \ \text{при} \ x \rightarrow x_0 \iff \exists \lim\limits_{x \to x_0} {\dfrac{\alpha(x)}{\beta(x)}} = c \in \mathbb{R} \backslash \{0\}
    $$
\item Функцию $\alpha(x)$ называют \textbf{бесконечно малой более высокого порядка} малости по сравнению с $\beta(x)$ при $x \rightarrow x_0$ и записывают $\alpha(x) = o(\beta(x))$, если существует и равен нулю предел отношения $\alpha(x) / \beta(x)$, при $x \rightarrow x_0$.

$$\alpha(x) = o(\beta(x)) \ \text{при} \ x \rightarrow x_0 \iff \exists \lim\limits_{x \to x_0} {\dfrac{\alpha(x)}{\beta(x)}} = 0
    $$
\item Функцию $\alpha(x)$ называют \textbf{бесконечно малой более низкого порядка} малости по сравнению с $\beta(x)$ при $x \rightarrow x_0$, если предел отношения $\alpha(x) / \beta(x)$, при $x \rightarrow x_0$, равен бесконечности.
\item Функции $\alpha(x)$ и $\beta(x)$ называют \textbf{эквивалентными} бесконечно малыми при $x \rightarrow x_0$, если предел их отношения при $x \rightarrow x_0$ равен 1.

$$\alpha(x) \sim \beta(x) \ \text{при} \ x \rightarrow x_0 \iff \lim\limits_{x \to x_0}{\dfrac{\alpha(x)}{\beta(x)}} = 1
    $$
\item Функцию $\alpha(x)$ называют \textbf{бесконечно малой $k$-ого порядка} малости относительно $\beta(x)$ при $x \rightarrow x_0$, а число $k \ (k > 0)$ - \textbf{порядком малости} $\alpha(x)$ относительно $\beta(x)$ при $x \rightarrow x_0$, если функции $\alpha(x)$ и $\beta^k(x)$ являются бесконечно малыми одного порядка при $x \rightarrow x_0$, т.е.

$$\lim\limits_{x \to x_0} {\dfrac{\alpha(x)}{\beta^k(x)}} = c \in \mathbb{R} \backslash \{0\}
    $$
\item Функцию $u(x)$ называют \textbf{бесконечно большой $k$-ого порядка} роста относительно $w(x)$ при $x \rightarrow x_0$, а число $k \ (k > 0)$ - \textbf{порядком роста} $u(x)$ относительно $w(x)$ при $x \rightarrow x_0$, если функции $u(x)$ и $w^k(x)$ являются бесконечно большими одного порядка при $x \rightarrow x_0$, т.е.

$$\lim\limits_{x \to x_0} {\dfrac{u(x)}{w^k(x)}} = c \in \mathbb{R} \backslash \{0\}
    $$
% правильно? 
\item \textbf{Главная часть суммы бесконечно малых функций} - это слагаемое более низкого порядка малости по сравнению с каждым из остальных слагаемых.
\item \textbf{Приращением аргумента} в точке $x_0$ называется изменение аргумента функции от значения $x_0$ к другому значению x,

$$\Delta x = x - x_0
    $$
\item \textbf{Приращением функции} в точке $x_0$ называется $\Delta y = f(x_0 + \Delta x) - f(x_0)$.
\item \textit{(опр. 1)} Функция $f(x)$ называется \textbf{непрерывной в точке $x_0$}, если в этой точке существует конечный предел функции и он совпадает с значением функции в этой точке, т.е. $\exists \lim\limits_{x \to x_0}f(x) = f(x_0)$.
\item \textit{(опр. 2)} Функция $f(x)$ называется \textbf{непрерывной в точке $x_0$}, если приращение функции в этой точке есть бесконечно малая функция при стремлении приращения аргумента к 0 ($\Delta x \rightarrow 0$).
\item Функция $f(x)$ называется \textbf{непрерывной в точке $x_0$ справа}, если в этой точке существует конечный \textit{правый} предел функции и он совпадает с значением функции в этой точке, т.е. $\exists \lim\limits_{x \to x_0+}f(x) = f(x_0)$.
\item Функция $f(x)$ называется \textbf{непрерывной в точке $x_0$ слева}, если в этой точке существует конечный левый предел функции и он совпадает с значением функции в этой точке, т.е. $\exists \lim\limits_{x \to x_0-}f(x) = f(x_0)$.
\item Функция $f(x)$ \textbf{непрерывна на интервале $(a, b)$}, если она непрерывна в каждой его точке.
\item Функция $f(x)$ \textbf{непрерывна на отрезке $[a, b]$}, если она непрерывна на интервале $(a, b)$, в точке $a$ - непрерывна справа, т.е. $\lim\limits_{x \to a+}f(x) = f(a)$, в точке $b$ - непрерывна слева, т.е. $\lim\limits_{x \to b-}f(x) = f(b)$.
\item Если данная функция $f(x)$ не является непрерывной в точке $x_0$, то $x_0$ называется \textbf{точкой разрыва} функции $f(x)$.
\item \textbf{Точкой разрыва первого рода} называют такую точку разрыва функции, в которой существуют оба односторонних предела этой функции и они конечны.
\item \textbf{Точкой разрыва второго рода} называют такую точку разрыва функции, в которой хотя бы один из односторонних пределов функции не существует (в частности, равен бесконечности).
\item Если $x_0$ — точка разрыва функции первого рода и односторонние пределы функции в этой точке равны между собой, но не равны значению функции в этой точке или $f(x)$ не определена в этой точке, то такой разрыв называют \textbf{устранимым}, а точку $x_0$ - \textbf{точкой устранимого разрыва первого рода}.
\item Если $x_0$ — точка разрыва функции первого рода и односторонние пределы функции в этой точке не равны между собой, то такой разрыв называют \textbf{неустранимым}, а точку $x_0$ - \textbf{точкой неустранимого разрыва первого рода}.
\item Если существует конечный предел $\lim\limits_{\Delta x \to 0} {\dfrac{f(x_0 + \Delta x) - f(x_0)}{\Delta x}}$, то он называется \textbf{производной функции $f(x)$ в точке} $x_0$ и обозначается $f'(x_0)$.
\item Если $f(x)$ определена в правосторонней окрестности точки $x_0$ и если

$\exists \lim\limits_{\Delta x \to 0+}{\dfrac{f(x_0+\Delta x) - f(x_0)}{\Delta x}}$, то этот предел называется \textbf{правой производной функции $f(x)$ в $x_0$} и обозначается $f'_+(x)$.
\item Если $f(x)$ определена в левосторонней окрестности точки $x_0$, и если

$\exists \lim\limits_{\Delta x \to 0-}{\dfrac{f(x_0+\Delta x) - f(x_0)}{\Delta x}}$, то этот предел называется \textbf{левой производной функции $f(x)$ в $x_0$} и обозначается $f'_-(x)$.
\item Пусть функция $y = f(x)$ определена в некоторой окрестности точки $x_0$. Функция называется \textbf{дифференцируемой в точке $x_0$}, если ее приращение $\Delta y$ в точке $x_0$ представимо в следующем виде: $\Delta y = A\Delta x + \alpha(\Delta x)\Delta x$, где $A$ - некоторое число, не зависящее от $\Delta x$, а $\lim\limits_{\Delta x \to 0}\alpha(\Delta x) = 0$.
\item \textbf{Дифференциалом функции $f(x)$ в точке $x_0$} называется главная часть приращения функции, линейная относительно приращения аргумента $\Delta x$.
\item \textbf{Дифференциалом $n$-го порядка} называется дифференциал от дифференциала $n-1$ порядка, т.е. $$d^{\ n}y = d(d^{\ n-1}y) = f^{(n)}(x)dx^n$$
\item \textbf{Производная $n$-ого порядка} от функции $y = f(x)$, есть производная от производной $n-1$ порядка, т.е. $$f^{(n)} = (f^{n-1}(x))'$$
\item Функция $f(x)$ называется \textbf{возрастающей на интервале} $(a, b)$, если $\forall x_1, x_2 \in (a, b)$, таких что $x_2 > x_1$, выполняется неравенство $f(x_2)>f(x_1)$.
\item Функция $f(x)$ называется \textbf{невозрастающей на интервале} $(a, b)$, если $\forall x_1, x_2 \in (a, b)$, таких что $x_2 > x_1$, выполняется неравенство $f(x_2)\leqslant f(x_1)$.
\item Функция $f(x)$ называется \textbf{убывающей на интервале} $(a, b)$, если $\forall x_1, x_2 \in (a, b)$, таких что $x_2 > x_1$, выполняется неравенство $f(x_2)<f(x_1)$.
\item Функция $f(x)$ называется \textbf{неубывающей на интервале} $(a, b)$, если $\forall x_1, x_2 \in (a, b)$, таких что $x_2 > x_1$, выполняется неравенство $f(x_2)\geqslant f(x_1)$.
\item Функция $f(x)$ называется \textbf{монотонной}, если она невозрастающая или неубывающая.
\item Функция $f(x)$ называется \textbf{строго монотонной}, если она возрастающая или убывающая.
\item Точка $x_0$ называется \textbf{точкой локального минимума} функции $f(x)$, если $\exists U_{\delta}(x_0)$, такая что $\forall x \in U_{\delta}(x_0): \ f(x_0) \leqslant f(x)$.
\item Точка $x_0$ называется \textbf{точкой локального максимума} функции $f(x)$, если $\exists U_{\delta}(x_0)$, такая что $\forall x \in U_{\delta}(x_0): \ f(x_0) \geqslant f(x)$.
\item Точка $x_0$ называется \textbf{точкой строгого локального минимума} функции $f(x)$, если $\exists \overset{\circ}{U}_{\delta}(x_0)$, такая что $\forall x \in \overset{\circ}{U}_{\delta}(x_0): \ f(x_0) < f(x)$.
\item Точка $x_0$ называется \textbf{точкой строгого локального максимума} функции $f(x)$, если $\exists \overset{\circ}{U}_\delta(x_0)$, такая что $\forall x \in \overset{\circ}{U}_\delta(x_0): \ f(x_0) > f(x)$.
\item \textbf{Точками локального экстремума} называются точки локального максимума и строгого локального максимума, локального минимума и строгого локального минимума.
\item \textbf{Точками строгого локального экстремума} называются точки строгого локального максимума и минимума.
\item Точку $x_0$ из области определения функции $f(x)$ называют \textbf{критической}, если производная в ней равна 0 или не сущестует вовсе.
\item Точку $x_0$ из области определения функции $f(x)$ называют \textbf{стационарной}, если $f'(x_0) = 0$.
\item Прямая $Ax + By + C = 0$ называется \textbf{асимптотой} графика $y = f(x)$, если расстояние от точки $M(x, f(x))$ графика функции до этой прямой стремится к 0 при бесконечном удалении точки $M$ от начала координат.
\item Прямая $x = a$ называется \textbf{вертикальной асимптотой} графика функции $y = f(x)$, если хотя бы один из пределов $\lim\limits_{x \to a+(-)}f(x) = \infty$
\item Прямая $y = kx + b$ называется \textbf{правой наклонной асимптотой} графика функции $y = f(x)$, если эту функцию можно представить в виде $f(x) = kx + b + \alpha(x)$, где $k, b \in \mathbb{R}$ и $\alpha(x)$ - бесконечно малая функция при $x \rightarrow +\infty$.
\item Прямая $y = kx + b$ называется \textbf{левой наклонной асимптотой} графика функции $y = f(x)$, если эту функцию можно представить в виде $f(x) = kx + b + \alpha(x)$, где $k, b \in \mathbb{R}$ и $\alpha(x)$ - бесконечно малая функция при $x \rightarrow -\infty$.
\item Пусть функция $f(x)$ дифференцируема на интервале $(a, b)$. График функции $y = f(x)$ имеет на интервале $(a, b)$ \textbf{выпуклость вверх}, если он лежит не выше любой касательной к графику на (a, b).
\item Пусть функция $f(x)$ дифференцируема на интервале $(a, b)$. График функции $y = f(x)$ имеет на интервале $(a, b)$ \textbf{выпуклость вниз}, если он лежит не ниже любой касательной к графику на (a, b).
\item Точка $x_0 \in (a,b)$ называется точкой перегиба функции $f(x)$, если эта функция непрерывна в точке $x_0$ и если $\exists \delta > 0$ такое, что направления выпуклостей функции $f(x)$ на интервалах $(x_0-\delta; x_0)$ и $(x_0; x_0+\delta)$ различны.

\end{enumerate}
\newpage 
\section*{Вопросы для подготовки к экзамену}\addcontentsline{toc}{section}{Вопросы для подготовки к экзамену}
\subsection*{1. \textit{Теорема (о единственности предела сходящейся последовательности)}}\addcontentsline{toc}{subsection}{1. \textit{Теорема (о единственности предела сходящейся последовательности)}}
\begin{Quote2} 
\small\centering 

Используются определения 15, 25, 26 из списка \end{Quote2} 

Если последовательность имеет предел, то этот предел - единственный.
\vspace*{20pt} 

\textit{Доказательство (от противного)}

Пусть $a, b \in \mathbb{R}, \ a \neq b$, где $a$ и $b$ - пределы сходящейся последовательности $\{X_n\}$:

$$
\lim\limits_{n \to \infty} x_n = a,\ \lim\limits_{n \to \infty} x_n = b, \ a \neq b
$$

По определению предела:

$$
\forall {\varepsilon} > 0, \ \exists N_1 = N_1({\varepsilon}) \in \mathbb{N} : \ \forall n > N_1 \Rightarrow  | x_n - a|  < {\varepsilon} 
$$

$$
\forall {\varepsilon} > 0, \ \exists N_2 = N_2({\varepsilon}) \in \mathbb{N} : \ \forall n > N_2 \Rightarrow  | x_n - b|  < {\varepsilon} 
$$

Примем ${\varepsilon} = {\dfrac{\mid  b - a\mid  }{3}}$ и при $n > max(N_1, \ N_2)$ получим

$$
| b - a|  = | x_n - a + b - x_n|  \leqslant | x_n - a|  + | b - x_n|  = | x_n - a|  + | x_n - b|  \Rightarrow | b - a|  < 2{\varepsilon}
$$

Или $\mid  b-a\mid   < 2 \cdot {\dfrac{\mid  b - a\mid  }{3}}$, т.е. $\mid  b-a\mid   < \dfrac{2}{3}\mid  b - a\mid  $, $\dfrac{1}{3}\mid  b-a\mid   < 0$, чего не может быть $\Rightarrow$ $a \neq b$ - неверно, т.е. $a = b \Rightarrow$ предел единственный.

Теорема доказана.
\newpage 
\subsection*{2. \textit{Теорема (об ограниченности сходящейся последовательности)}}\addcontentsline{toc}{subsection}{2. \textit{Теорема (об ограниченности сходящейся последовательности)}}
\begin{Quote2} 
\small\centering 

Используются определения 15, 24, 25, 26 из списка \end{Quote2} 

Всякая сходящаяся последовательность является ограниченной.
\vspace*{20pt} 

\textit{Доказательство}

Пусть $\{X_n\}$ - сходящаяся последовательность. Тогда по определению, у нее существует конечный предел

$$
\lim\limits_{n \to \infty}x_n = a \iff \forall {\varepsilon} > 0 \ \exists N = N({\varepsilon})\in \mathbb{N}: \ \forall n > \mathbb{N} \Rightarrow | x_n - a|  < {\varepsilon},
$$

$$
-{\varepsilon}+a < x_n < {\varepsilon} + a
$$

Обозначим через $A$ максимальное число среди $\mid  x_1\mid  , \mid  x_2\mid  , ..., \mid  x_n\mid  , \mid  a-{\varepsilon}\mid  , \mid  a + {\varepsilon}\mid  $, т.е.

$$
A = max(| x_1| , | x_2| , ..., | x_n| , | a-{\varepsilon}| , | a + {\varepsilon}| )
$$

Тогда $\forall n \in \mathbb{N}$ выполняется $\mid  x_n\mid   < A, \Rightarrow$ последовательность ограничена.

Теорема доказана.
\newpage 
\subsection*{3. \textit{Теорема (о локальной ограниченности функции, имеющей конечный предел)}}\addcontentsline{toc}{subsection}{3. \textit{Теорема (о локальной ограниченности функции, имеющей конечный предел)}}
\begin{Quote2} 
\small\centering 

Используются определения 28, 34 из списка \end{Quote2} 

Если функция $f(x)$ имеет конечный предел при $x \rightarrow x_0$, то $f(x)$ локально ограничена.
\vspace*{20pt} 

\textit{Доказательство}

По условию $\exists$ конечный предел $\lim\limits_{x \to x_0}f(x) = a$, тогда

$$
\forall {\varepsilon} > 0 \ \exists\delta=\delta({\varepsilon} )>0: \ \forall x \in \overset{\circ}U_\delta(x_0) \Rightarrow | f(x) - a|  < {\varepsilon}
$$

Пусть ${\varepsilon} = 1$, тогда $\mid  f(x)\mid   - \mid  a\mid   \leqslant \mid  f(x) - a\mid   < 1$, а значит

$$
\forall x \in \overset{\circ}U_\delta(x_0) \Rightarrow | f(x)|  <1 + | A|  = const \overset{\text{по опр.}}{\Rightarrow}$$

$f(x)$ является локально ограниченной в окрестности точки $x_0$.

Теорема доказана.
\newpage 
\subsection*{4. \textit{Теорема (о сохранении функцией знака своего предела)}}\addcontentsline{toc}{subsection}{4. \textit{Теорема (о сохранении функцией знака своего предела)}}
\begin{Quote2} 
\small\centering 

Используются определения 28 из списка \end{Quote2} 

Если $\lim\limits_{x \to x_0} f(x) = A \neq 0$, то $\exists \overset{\circ}U_\delta(x_0): \ \forall x \in \overset{\circ}U_\delta(x_0)$ функция $f(x)$ сохраняет знак своего предела.
\vspace*{20pt} 

\textit{Доказательство}

По условию $\exists$ конечный $\lim\limits_{x \to x_0} f(x) = a > 0 \overset{\text{по опр.}}{\Rightarrow}$

$$
\forall {\varepsilon} > 0 \ \exists\delta=\delta({\varepsilon} )>0: \ \forall x \in \overset{\circ}U_\delta(x_0) \Rightarrow | f(x) - a|  < {\varepsilon}
$$
\begin{itemize}

\item в случае $a > 0$ выбираем ${\varepsilon} = {\dfrac{a}{2}}$, тогда

$$| f(x) - a|  < {\dfrac{a}{2}}
    $$$$-{\dfrac{a}{2}}< f(x) - a < {\dfrac{a}{2}}
    $$$${\dfrac{a}{2}}< f(x) < {\dfrac{3a}{2}}
    $$

Следовательно $f(x) > {\dfrac{a}{2}}>0$, т.е. данная функция положительна при $x \in \overset{\circ}U_\delta(x_0)$.
\item в случае $a < 0$ выбираем ${\varepsilon} = -{\dfrac{a}{2}}$, тогда

$$| f(x) - a|  < - {\dfrac{a}{2}}
    $$$${\dfrac{a}{2}}< f(x) - a < -{\dfrac{a}{2}}
    $$$${\dfrac{3a}{2}}< f(x) < {\dfrac{a}{2}}
    $$

Следовательно $f(x) < {\dfrac{a}{2}}< 0$, т.е. данная функция отрицательна при $x \in \overset{\circ}U_\delta(x_0)$.

\end{itemize}

Теорема доказана.
\newpage 
\subsection*{5. \textit{Теорема (о предельном переходе в неравенстве)}}\addcontentsline{toc}{subsection}{5. \textit{Теорема (о предельном переходе в неравенстве)}}
\begin{Quote2} 
\small\centering 

Используются определения 28 из списка \end{Quote2} 

Пусть функции $f(x)$ и $g(x)$ определены в проколотой окрестности $\overset{\circ}U(x_0)$ точки $x_0$, причем для любого $x \in \overset{\circ}U(x_0)$ выполняется неравенство $f(x) \geqslant g(x)$. Тогда, если эти функции имеют пределы $a = \lim\limits_{x \to x_0}f(x)$ и $b = \lim\limits_{x \to x_0}g(x)$, то $a \geqslant b$.
\vspace*{20pt} 

\textit{Доказательство}

По условию $\forall x \in \overset{\circ}U(x_0): \ f(x) \geqslant g(x) \Rightarrow f(x) - g(x) \geqslant 0$, тогда по теореме о сохранении функцией знака своего предела:

$$
\lim\limits_{x \to x_0}(f(x) - g(x)) \geqslant 0 \Rightarrow \lim\limits_{x \to x_0}f(x) - \lim\limits_{x \to x_0}g(x) = a - b \geqslant 0, \Rightarrow a \geqslant b
$$

Теорема доказана.
\newpage 
\subsection*{6. \textit{Теорема (о пределе промежуточной функции)}}\addcontentsline{toc}{subsection}{6. \textit{Теорема (о пределе промежуточной функции)}}
\begin{Quote2} 
\small\centering 

Используются определения 28 из списка \end{Quote2} 

Пусть для всех $x$ из некоторой проколотой окрестности $\overset{\circ}U(x_0)$ точки $x_0$ выполняется двойное неравенство $f(x) \leqslant g(x) \leqslant h(x)$, и пусть существуют пределы $\lim\limits_{x \to x_0}f(x)$ и $\lim\limits_{x \to x_0}h(x)$, равные одному и тому же числу $a$. Тогда и $\lim\limits_{x \to x_0}g(x) = a$.
\vspace*{20pt} 

\textit{Доказательство}

По условию $\exists \lim\limits_{x \to x_0}f(x) = a$, $\lim\limits_{x \to x_0}h(x) = a$, тогда по определению предела функции,

$$
\forall {\varepsilon}  > 0 \ \exists \delta_1= \delta_1 ({\varepsilon} ) > 0: \ \forall x \in \overset{\circ}U_{\delta_1}(x_0) \Rightarrow | f(x) - a|  < {\varepsilon}
$$

$$
\text{т.е.} \ \ a - {\varepsilon} < f(x) < a + {\varepsilon}
$$

$$
\forall {\varepsilon}  > 0 \ \exists \delta_2= \delta_2 ({\varepsilon} ) > 0: \ \forall x \in \overset{\circ}U_{\delta_2}(x_0) \Rightarrow | h(x) - a|  < {\varepsilon}
$$

$$
\text{т.е.} \ \ a - {\varepsilon} < h(x) < a + {\varepsilon}
$$

Тогда при $x \in \overset{\circ}U_{\delta}(x_0), \ \delta = min(\delta_1, \delta_2)$, выполняется неравенство

$$
a - {\varepsilon} < f(x) \leqslant g(x) \leqslant h(x) < a+ {\varepsilon}
$$ $$
a - {\varepsilon} <  g(x) < a+ {\varepsilon}
$$ $$
| g(x) - a|  < {\varepsilon}
$$

Таким образом, получаем

$$
\forall {\varepsilon} > 0 \ \exists \delta = \delta({\varepsilon}) > 0: \ \forall x \in \overset{\circ}U_{\delta}(x_0) \Rightarrow | g(x) - a|  < {\varepsilon} \iff \lim\limits_{x \to x_0} g(x) = a
$$

Теорема доказана.
\newpage 
\subsection*{7. \textit{Теорема (о пределе произведения функций)}}\addcontentsline{toc}{subsection}{7. \textit{Теорема (о пределе произведения функций)}}
\begin{Quote2} 
\small\centering 

Используются определения 28, 35 из списка \end{Quote2} 

Если $\exists$ конечные пределы $\lim\limits_{x \to x_0}f(x)=a$ и $\lim\limits_{x \to x_0}g(x) = b$, то $\lim\limits_{x \to x_0}(f(x)\cdot g(x)) = a\cdot b = \lim\limits_{x \to x_0}f(x)\cdot \lim\limits_{x \to x_0}g(x)$.
\vspace*{20pt} 

\textit{Доказательство}

По условию $\exists$ конечные пределы $\lim\limits_{x \to x_0}f(x)=a$ и $\lim\limits_{x \to x_0}g(x) = b$, тогда по теореме о связи функции, ее предела и бесконечно малой имеем

$$
f(x) = a + \alpha(x), \ \text{где} \ \alpha(x) - \text{бесконечно малая функция при} \ x \rightarrow x_0 
$$

$$
g(x) = b + \beta(x), \ \text{где} \ \beta(x) - \text{бесконечно малая функция при} \ x \rightarrow x_0 
$$

Тогда

$$
f(x)\cdot g(x) = (a + \alpha(x))\cdot(b + \beta(x)) = a\cdot b + a\cdot \beta(x) + \alpha(x)\cdot b + \alpha(x)\cdot \beta(x)
$$

$$
\lim\limits_{x \to x_0}(f(x)\cdot g(x)) = \lim\limits_{x \to x_0}(a\cdot b + \underbrace{a\cdot \beta(x)}_{\text{б.м. при} \ x \rightarrow x_0} + \underbrace{\alpha(x)\cdot b}_{\text{б.м. при} \ x \rightarrow x_0} + \underbrace{\alpha(x)\cdot \beta(x)}_{\text{б.м. при} \ x \rightarrow x_0}) =
$$

$$
= \lim\limits_{x \to x_0}(a\cdot b) = a\cdot b = \lim\limits_{x \to x_0}f(x) \cdot \lim\limits_{x \to x_0}g(x)
$$

Теорема доказана.
\newpage 
\subsection*{8. \textit{Теорема (о пределе сложной функции)}}\addcontentsline{toc}{subsection}{8. \textit{Теорема (о пределе сложной функции)}}
\begin{Quote2} 
\small\centering 

Используются определения 25, 29 из списка \end{Quote2} 

Если функция $y = f(x)$ имеет в точке $x = a$ конечный предел, равный $b$, и $f(x) \neq b$ в некоторой проколотой окрестности $\overset{\circ}U(a)$ этой точки, а функция $g(y)$ имеет в точке $b$ конечный предел $c$, то сложная функция $g(f(x))$ имеет $\lim\limits_{x \to a}g(f(x)) = c$.
\vspace*{20pt} 

\textit{Доказательство}

По определению предела функции по Гейне имеем:

$$
\exists \lim\limits_{x \to a}f(x) = b \iff \{\forall x_n \in \overset{\circ}U(a), \ n \in \mathbb{N}\} \ \cap \ \lim\limits_{n \to \infty}{x_n} = a : \ \lim\limits_{n \to \infty}\{f(x_n)\} = b
$$

$$
\exists \lim\limits_{y \to b}g(y) = c \iff \{\forall y_n \in \overset{\circ}U(b), \ n \in \mathbb{N}\} \ \cap \ \lim\limits_{n \to \infty}{y_n} = b : \ \lim\limits_{n \to \infty}\{g(y_n)\} = c
$$

Пусть $\{X_n\}$ - произвольная последовательность, стремящаяся к точке $a$ и $x_n \neq a \ \forall n \in \mathbb{N}$. Тогда $\lim\limits_{n \to \infty}{\{f(x_n)\}} = b$, но $f(x_n) \neq b \ \forall n \in \mathbb{N}$. Пусть $y_n = f(x_n)$. Поскольку $\lim\limits_{n \to \infty}{\{y_n\}} = b$ и $y_n \neq b \ \forall n \in \mathbb{N}$, имеем $\lim\limits_{n \to \infty}{\{g(y_n)\}} = c$, т.е.

$$
\{\forall x_n \in \overset{\circ}U(a), \ n \in \mathbb{N}\} \ \cap \ \lim\limits_{n \to \infty}{x_n} = a : \ \lim\limits_{n \to \infty}\{g(f(x_n))\} = c \iff \lim\limits_{n \to \infty} g(f(x)) = c
$$

Теорема доказана.
\newpage 
\subsection*{9. \textit{Вывод первого замечательного предела}}\addcontentsline{toc}{subsection}{9. \textit{Вывод первого замечательного предела}}
\begin{mainQuote}

$$ \lim\limits_{x \to 0} {\dfrac{sin(x)}{x}} = 1 $$
\end{mainQuote}

Пусть $0 < x < {\dfrac{\pi}{2}}$. Рассмотрим окружность радиуса $R$ с центром в начале координат, пересекающую ось абсцисс в точке $A$, и пусть угол $\angle AOB$ равен $x$ (радиан). Пусть, далее, $CA$ — перпендикуляр к этой оси, $C$ — точка пересечения с этим перпендикуляром продолжения отрезка $OB$ за точку $B$. Тогда

$$S_{\triangle OAB} < S_{сектор \ OAB} < S_{\triangle OAC}$$ $$\dfrac{1}{2}R^2sin(x) <\dfrac{1}{2}R^2x < \dfrac{1}{2}R^2tg(x)$$ $$sin(x) < x < tg(x)$$ $$1 < \dfrac{x}{sin(x)} < \dfrac{1}{cos(x)}$$ $$1 > \dfrac{sin(x)}{x} > cos(x), \ \text{при} \ x \in (0, \dfrac{\pi}{2})$$

Рассмотрим $x \in (-{\dfrac{\pi}{2}}, 0)$. Сделаем замену $\beta = -x$, таким образом $\beta \in (0, {\dfrac{\pi}{2})}$, а значит справедливо следующее неравенство: $$1 > \dfrac{sin(\beta)}{\beta} > cos(\beta)$$ Вернемся к замене $\beta = -x$

$$1 > \dfrac{sin(-x)}{-x} > cos(-x)$$ $$1 > \dfrac{-sin(x)}{-x} > cos(x)$$ $$1 > \dfrac{sin(x)}{x} > cos(x) \ \text{при} \ x \in (-{\dfrac{\pi}{2}}, 0)$$

Таким образом, полученное неравенство справедливо для $x \in ({-{\dfrac{\pi}{2}}}, 0)\cup(0, {\dfrac{\pi}{2}})$. Перейдем к пределу при $x \rightarrow 0$:

$$
\begin{rcases}\lim\limits_{x \to 0}cos(x) = 1\\ 
\lim\limits_{x \to 0}1 = 1
\end{rcases} \Rightarrow \ \text{(по т. о пределе промежуточной функции)} \ \lim\limits_{x \to 0}\dfrac{sin(x)}{x} = 1 
$$
\newpage 
\subsection*{10. \textit{Теорема (о связи функции, ее предела и бесконечно малой)}}\addcontentsline{toc}{subsection}{10. \textit{Теорема (о связи функции, ее предела и бесконечно малой)}}
\begin{Quote2} 
\small\centering 

Используются определения 28, 35 из списка \end{Quote2} 

Равенство $\lim\limits_{x \to x_0}f(x) = a$ имеет место $\iff f(x) = a + \alpha(x)$, где $\alpha(x)$ - бесконечно малая функция при $x \rightarrow x_0$.
\vspace*{20pt} 

\textit{Доказательство}

$(\Rightarrow)$

По условию $\exists$ конечный $\lim\limits_{x \to x_0}f(x) = a$, тогда по определению

$$
\forall{\varepsilon} > 0 \ \exists \overset{\circ}U_\delta(x_0): \ \forall x \in \overset{\circ}U_\delta(x_0) \Rightarrow |f(x) - a| <{\varepsilon}
$$

Обозначим $f(x) - a = \alpha(x)$. Тогда $\mid \alpha(x)\mid  < {\varepsilon}\ \ \forall x \in \overset{\circ}U_\delta(x_0) \overset{\text{по опр.}}\Rightarrow \lim\limits_{x \to x_0}\alpha(x) = 0,$ т.е. $\alpha(x)$ - бесконечно малая функция при $x \rightarrow x_0$.

Но $\alpha(x) = f(x) - a \Rightarrow f(x) = a + \alpha(x), \text{ где } \lim\limits_{x \to x_0}\alpha(x) = 0$.

$(\Leftarrow)$

По условию $f(x) = a + \alpha(x),\text{ где } \lim\limits_{x \to x_0}\alpha(x) = 0$, тогда по определению

$$
\forall{\varepsilon} > 0 \ \exists \overset{\circ}U_\delta(x_0): \ \forall x \in \overset{\circ}U_\delta(x_0) \Rightarrow |\alpha(x)| <{\varepsilon}
$$

Но по условию $f(x) = a + \alpha(x) \Rightarrow \alpha(x) = f(x) - a$, отсюда имеем

$$
\forall{\varepsilon} > 0 \ \exists \overset{\circ}U_\delta(x_0): \ \forall x \in \overset{\circ}U_\delta(x_0) \Rightarrow |f(x) - a| <{\varepsilon} \overset{\text{по опр.}}\Rightarrow \lim\limits_{x \to x_0}f(x) = a
$$

Теорема доказана.
\newpage 
\subsection*{11. \textit{Теорема (о произведении бесконечно малой функции на ограниченную)}}\addcontentsline{toc}{subsection}{11. \textit{Теорема (о произведении бесконечно малой функции на ограниченную)}}
\begin{Quote2} 
\small\centering 

Используются определения 34, 35 из списка \end{Quote2} 

Если $\alpha(x)$ - бесконечно малая функция при $x \rightarrow x_0$, $f(x)$ - ограниченная функция, то $\alpha(x)\cdot f(x)$ - бесконечно малая функция при $x \rightarrow x_0$.
\vspace*{20pt} 

\textit{Доказательство}

По условию $\alpha(x)$ - бесконечно малая функция при $x \rightarrow x_0$, тогда

$$
\forall {\varepsilon} > 0 \ \exists\overset{\circ}U_1(x_0): 
\forall x \in \overset{\circ}U_1(x_0) \Rightarrow |\alpha(x)| < \dfrac{ {\varepsilon}}{c}
$$

$f(x)$ - ограниченная функция, тогда

$$
|f(x)| < c, \text{ где } c = const, \ \forall x \in \overset{\circ}U_2(x_0)
$$

Таким образом,

$$
\forall x \in \overset{\circ}U(x_0) = \overset{\circ}U_1(x_0) \cap \overset{\circ}U_2(x_0):
$$ $$
|\alpha(x)\cdot f(x)| < \dfrac{ {\varepsilon}}{c} \cdot c \Rightarrow
$$

$ \mid \alpha(x)\cdot f(x)\mid  < {\varepsilon} \Rightarrow \alpha(x)\cdot f(x)$ - бесконечно малая функция при $x \rightarrow x_0$

Теорема доказана.
\newpage 
\subsection*{12. \textit{Теорема (о связи между бесконечно большой и бесконечно малой)}}\addcontentsline{toc}{subsection}{12. \textit{Теорема (о связи между бесконечно большой и бесконечно малой)}}
\begin{Quote2} 
\small\centering 

Используются определения 28, 35, 36 из списка \end{Quote2} 

$\alpha(x)$ - бесконечно малая функция при $x \rightarrow x_0$, отличная от нуля в некоторой проколотой окрестности точки $x_0 \ \Rightarrow \ {\dfrac{1}{\alpha(x)}}$ - бесконечно большая функция при $x \rightarrow x_0$.

$f(x)$ - бесконечно большая функция при $x \rightarrow x_0 \ \Rightarrow \ {\dfrac{1}{f(x)}}$ - бесконечно малая функция при $x \rightarrow x_0$.
\vspace*{20pt} 

\textit{Доказательство}

Пусть $\alpha(x)$ - бесконечно малая функция при $x \rightarrow x_0$, отличная от нуля в некоторой проколотой окрестности $\overset{\circ}U(x_0)$ точки $x_0$. Выберем произвольное $E > 0$. Тогда

$$
\text{для } {\varepsilon} = \dfrac{1}{E} > 0 \ \ \exists \overset{\circ}U_1(x_0): \ \forall x \in \overset{\circ}U(x_0) \cap \overset{\circ}U_1(x_0) \Rightarrow 0 < |\alpha(x)| < {\varepsilon}, \text{ т.е.}
$$ $$
\dfrac{1}{|\alpha(x)|} > E, \overset{\text{по опр.}}\Rightarrow \dfrac{1}{\alpha(x)} \text{ - бесконечно большая функция при } x \rightarrow x_0 
$$

Пусть $f(x)$ - бесконечно большая функция при $x \rightarrow x_0$. Выберем произвольное ${\varepsilon} > 0$. Тогда

$$
\text{для } E = \dfrac{1}{{\varepsilon}} \ \ \exists \overset{\circ}U(x_0): \ \forall x \in \overset{\circ}U(x_0) \Rightarrow |f(x)| > E, \text{ т.е.}
$$ $$
\dfrac{1}{f(x)} < \dfrac{1}{E} = {\varepsilon}, \overset{\text{по опр.}}\Rightarrow \dfrac{1}{f(x)} \text{ - бесконечно малая функция при } x \rightarrow x_0 
$$

Теорема доказана.
\newpage 
\subsection*{13. \textit{Теорема (о замене бесконечно малой на эквивалентную под знаком предела)}}\addcontentsline{toc}{subsection}{13. \textit{Теорема (о замене бесконечно малой на эквивалентную под знаком предела)}}
\begin{Quote2} 
\small\centering 

Используются определения 35, 44 из списка \end{Quote2} 

Пусть $\alpha(x)\sim\beta(x)$ при $x \rightarrow x_0$, и $f(x)$ - некоторая функция, определенная в проколотой окрестности $\overset{\circ}U(x_0)$ точки $x_0$. Тогда:
\begin{itemize}

\item если существует предел при $x \rightarrow x_0$ произведения $\alpha(x)\cdot f(x)$, то он не изменится при замене $\alpha(x)$ на эквивалентную при $x \rightarrow x_0$ бесконечно малую функцию $\beta(x)$
\item если существует предел при $x \rightarrow x_0$ частного ${\dfrac{f(x)}{\alpha(x)}}$ , то он не изменится при замене $\alpha(x)$ на эквивалентную при $x \rightarrow x_0$ бесконечно малую функцию $\beta(x)$

\end{itemize}
\vspace*{20pt} 

\textit{Доказательство}

По условию $\alpha(x)\sim\beta(x)$ при $x \rightarrow x_0$, тогда по определению $\lim\limits_{x \to x_0}{\dfrac{\alpha(x)}{\beta(x)}} = 1$. Таким образом,
\begin{itemize}

\item $\lim\limits_{x \to x_0}(\alpha(x)\cdot f(x)) = \lim\limits_{x \to x_0}{\dfrac{\alpha(x)\cdot\beta(x)\cdot f(x)}{\beta(x)}} = \lim\limits_{x \to x_0}{\dfrac{\alpha(x)}{\beta(x)}}\cdot\lim\limits_{x \to x_0}(\beta(x)\cdot f(x)) = \lim\limits_{x \to x_0}(\beta(x)\cdot f(x))$
\item $\lim\limits_{x \to x_0}{\dfrac{f(x)}{\alpha(x)}} = \lim\limits_{x \to x_0}{\dfrac{\beta(x)\cdot f(x)}{\alpha(x)\cdot\beta(x)}} = \lim\limits_{x \to x_0}{\dfrac{\beta(x)}{\alpha(x)}}\cdot\lim\limits_{x \to x_0}{\dfrac{f(x)}{\beta(x)}} = \lim\limits_{x \to x_0}{\dfrac{f(x)}{\beta(x)}}$.

\end{itemize}

Теорема доказана.
\newpage 
\subsection*{14. \textit{Теорема (о необходимом и достаточном условии эквивалентности бесконечно малых)}}\addcontentsline{toc}{subsection}{14. \textit{Теорема (о необходимом и достаточном условии эквивалентности бесконечно малых)}}
\begin{Quote2} 
\small\centering 

Используются определения 35, 42, 44 из списка \end{Quote2} 

Две бесконечно малые функции $\alpha(x)$ и $\beta(x)$ при $x \rightarrow x_0$ эквивалентны $\iff$ их разность имеет больший порядок малости при $x \rightarrow x_0$ по сравнению с каждой из них.
\vspace*{20pt} 

\textit{Доказательство}

$(\Rightarrow)$

По условию $\alpha(x) \sim \beta(x)$ при $x \rightarrow x_0$, тогда по определению $\lim\limits_{x \to x_0}{\dfrac{\alpha(x)}{\beta(x)}} = 1$. Таким образом,

$$
\lim\limits_{x \to x_0}{\dfrac{\alpha(x) - \beta(x)}{\beta(x)}} = \lim\limits_{x \to x_0}{\dfrac{\alpha(x)}{\beta(x)}} - 1 = 0, \overset{\text{по опр.}}\Rightarrow \alpha(x) - \beta(x) = o (\beta(x)) \text{ при } x \rightarrow x_0
$$

$$
\lim\limits_{x \to x_0}{\dfrac{\alpha(x) - \beta(x)}{\alpha(x)}} = 1 - \lim\limits_{x \to x_0}{\dfrac{\beta(x)}{\alpha(x)}} = 0, \overset{\text{по опр.}}\Rightarrow \alpha(x) - \beta(x) = o (\alpha(x)) \text{ при } x \rightarrow x_0
$$

$(\Leftarrow)$

По условию $\alpha(x) - \beta(x) = o (\beta(x)) \text{ при } x \rightarrow x_0, \ \alpha(x) - \beta(x) = o (\alpha(x)) \text{ при } x \rightarrow x_0$. Тогда

$$
0 = \lim\limits_{x \to x_0}{\dfrac{\alpha(x) - \beta(x)}{\beta(x)}} = \lim\limits_{x \to x_0}{\dfrac{\alpha(x)}{\beta(x)}} - 1 \ \Rightarrow \lim\limits_{x \to x_0}{\dfrac{\alpha(x)}{\beta(x)}} = 1, \overset{\text{по опр.}}\Rightarrow \alpha(x) \sim \beta(x) \text{ при } x \rightarrow x_0
$$

$$
0 = \lim\limits_{x \to x_0}{\dfrac{\alpha(x) - \beta(x)}{\alpha(x)}} = 1 - \lim\limits_{x \to x_0}{\dfrac{\beta(x)}{\alpha(x)}} \ \Rightarrow \lim\limits_{x \to x_0}{\dfrac{\beta(x)}{\alpha(x)}} = 1, \overset{\text{по опр.}}\Rightarrow \alpha(x) \sim \beta(x) \text{ при } x \rightarrow x_0
$$

Теорема доказана.
\newpage 
\subsection*{15. \textit{Теорема (о сумме конечного числа бесконечно малых разных порядков)}}\addcontentsline{toc}{subsection}{15. \textit{Теорема (о сумме конечного числа бесконечно малых разных порядков)}}
\begin{Quote2} 
\small\centering 

Используются определения 35, 44 из списка \end{Quote2} 

Сумма конечного числа бесконечно малых функций при $x \rightarrow x_0$ эквивалентна своей главной части.
\vspace*{20pt} 

\textit{Доказательство}

Пусть $\alpha_1(x), \ \alpha_2(x), \ ..., \ \alpha_n(x)$ - бесконечно малые функции при $x \rightarrow x_0$, и $\alpha_1(x)$ - главная часть суммы $\alpha_1(x) + \alpha_2(x) + ... +  \alpha_n(x)$, т.е.

$$
\lim\limits_{x \to x_0}\dfrac{\alpha_2(x)}{\alpha_1(x)} = 0, \lim\limits_{x \to x_0}\dfrac{\alpha_3(x)}{\alpha_1(x)} = 0, \ ...,\ \lim\limits_{x \to x_0}\dfrac{\alpha_n(x)}{\alpha_1(x)} = 0, 
$$

Тогда рассмотрим

$$
\lim\limits_{x \to x_0}{\dfrac{\alpha_1(x) + \alpha_2(x) + ... +  \alpha_n(x)}{\alpha_1(x)}}= 1 + \lim\limits_{x \to x_0}\dfrac{\alpha_2(x)}{\alpha_1(x)} + \lim\limits_{x \to x_0}\dfrac{\alpha_3(x)}{\alpha_1(x)}+ ... + \lim\limits_{x \to x_0}\dfrac{\alpha_n(x)}{\alpha_1(x)} = 1, \overset{\text{по опр.}}\Rightarrow
$$

$$
\alpha_1(x) + \alpha_2(x) + ... +  \alpha_n(x) \sim \alpha_1(x) \text{ при } x \rightarrow x_0
$$

Теорема доказана.
\newpage 
\subsection*{16. \textit{Теорема (о непрерывности суммы, произведения и частного непрерывных функций)}}\addcontentsline{toc}{subsection}{16. \textit{Теорема (о непрерывности суммы, произведения и частного непрерывных функций)}}
\begin{Quote2} 
\small\centering 

Используются определения 50 из списка \end{Quote2} 

Если $f(x)$ и $g(x)$ непрерывны в точке $x_0$, то функции $f(x) \pm g(x),\ f(x)\cdot g(x), \ {\dfrac{f(x)}{g(x)}}$ (последнее при $g(x) \neq 0$) - также непрерывны в точке $x_0$.
\vspace*{20pt} 

\textit{Доказательство}

По условию $f(x)$ и $g(x)$ непрерывны в точке $x_0$, тогда по определению

$$\exists\lim\limits_{x \to x_0}f(x) = f(x_0)$$ $$\exists\lim\limits_{x \to x_0}g(x) = g(x_0)$$
\begin{enumerate}

\item $\lim\limits_{x \to x_0}(f(x) \pm g(x)) = \lim\limits_{x \to x_0}f(x) \pm \lim\limits_{x \to x_0}g(x) = f(x_0) \pm g(x_0), \overset{\text{по опр.}}\Rightarrow f(x) \pm g(x)$ - непрерывна в точке x = $x_0$.
\item $\lim\limits_{x \to x_0}(f(x) \cdot g(x)) = \lim\limits_{x \to x_0}f(x) \cdot \lim\limits_{x \to x_0}g(x) = f(x_0) \cdot g(x_0), \overset{\text{по опр.}}\Rightarrow f(x) \cdot g(x)$ - непрерывна в точке x = $x_0$.
\item $\lim\limits_{x \to x_0}{\dfrac{f(x)}{g(x)}} = {\dfrac{\lim\limits_{x \to x_0}f(x)} {\lim\limits_{x \to x_0}g(x)}} = {\dfrac{f(x_0)}{g(x_0)}}, \overset{\text{по опр.}}\Rightarrow {\dfrac{f(x)}{g(x)}}$ - непрерывна в точке x = $x_0$ при условии $g(x) \neq 0$.

\end{enumerate}

Теорема доказана.
\newpage 
\subsection*{17. \textit{Теорема (о непрерывности сложной функции)}}\addcontentsline{toc}{subsection}{17. \textit{Теорема (о непрерывности сложной функции)}}
\begin{Quote2} 
\small\centering 

Используются определения 50 из списка \end{Quote2} 

Если функция $y = f(x)$ непрерывна в точке $x = a$, а функция $g(y)$ непрерывна в соответствующей точке $b = f(a)$, то сложная функция $g(f(x))$ непрерывна в точке $x = a$.
\vspace*{20pt} 

\textit{Доказательство}

По условию функция $y = f(x)$ непрерывна в точке $x = a$, функция $g(y)$ непрерывна в точке $b = f(a)$. Тогда по определению $$\exists\lim\limits_{x \to a}f(x) = f(a)$$ $$\exists\lim\limits_{y \to b}g(y) = g(b)$$

Тогда

$$
\lim\limits_{x \to a}g(f(x)) = \lim\limits_{y \to b}g(y) = g(b) = g(f(a)), \overset{\text{по опр.}}\Rightarrow \text{функция } g(f(x)) \text{ непрерывна в точке } x = a
$$

Теорема доказана.
% так как проколотая окр = окр \ {x\_0}... 
\newpage 
\subsection*{18. \textit{Теорема (о сохранении знака непрерывной функции в окрестности точки)}}\addcontentsline{toc}{subsection}{18. \textit{Теорема (о сохранении знака непрерывной функции в окрестности точки)}}
\begin{Quote2} 
\small\centering 

Используются определения 50 из списка \end{Quote2} 

Пусть функция $f(x)$ непрерывна в точке $x_0$, и $f(x_0) \neq 0$. Тогда в некоторой окрестности $U_\delta(x_0)$ точки $x_0$ функция $f(x)$ имеет знак числа $f(x_0)$.
\vspace*{20pt} 

\textit{Доказательство}

По условию $y = f(x)$ непрерывна в точке $x = x_0$. Тогда по определению

$$
\exists\lim\limits_{x \to x_0}f(x) = f(x_0) \neq 0
$$

Тогда по теореме о сохранении функцией знака своего предела функция $f(x)$ имеет знак числа $f(x_0)$ в некоторой проколотой окрестности $\overset{\circ}U_\delta(x_0)$ точки $x_0$, т.е.

$$
f(x_0) > 0 \Rightarrow \exists \overset{\circ}U_\delta(x_0):\ \forall x \in \overset{\circ}U_\delta(x_0)\Rightarrow f(x) > 0 
$$

$$
f(x_0) < 0 \Rightarrow \exists \overset{\circ}U_\delta(x_0):\ \forall x \in \overset{\circ}U_\delta(x_0)\Rightarrow f(x) < 0 
$$

Так как $\overset{\circ}U_\delta(x_0) = U_\delta(x_0) \backslash \{x_0\}$, то

$$
f(x_0) > 0 \Rightarrow \exists U_\delta(x_0):\ \forall x \in U_\delta(x_0)\Rightarrow f(x) > 0 
$$ $$
f(x_0) < 0 \Rightarrow \exists U_\delta(x_0):\ \forall x \in U_\delta(x_0)\Rightarrow f(x) < 0 
$$

Теорема доказана.
\newpage 
\subsection*{19. \textit{Функция, непрерывная в точке. Теорема о непрерывности элементарных функций.}}\addcontentsline{toc}{subsection}{19. \textit{Функция, непрерывная в точке. Теорема о непрерывности элементарных функций.}}

\textit{(опр. 1)} Функция $f(x)$ называется \textbf{непрерывной в точке $x_0$}, если в этой точке существует конечный предел функции и он совпадает с значением функции в этой точке, т.е. $\exists \lim\limits_{x \to x_0}f(x) = f(x_0)$.

\textit{(опр. 2)} Функция $f(x)$ называется \textbf{непрерывной в точке $x_0$}, если приращение функции в этой точке есть бесконечно малая функция при стремлении приращения аргумента к 0 ($\Delta x \rightarrow 0$).

\textit{Теорема (о непрерывности элементарных функций)}

Все элементарные функции непрерывны всюду, где они определены.
\vspace*{20pt} 

\textit{Доказательство $(\ $для $y = sin(x)$ и $y = cos( x)\ )$}
\vspace*{20pt} 

Рассмотрим $y = sin(x)$:

Найдем приращение функции

$$
\Delta y = f(x + \Delta x) - f(x) = sin(x + \Delta x) - sin(x) =
$$ $$
= 2sin{\Big(}\dfrac{x + \Delta x - x}{2}{\Big)}\cdot cos\Big(\dfrac{x + \Delta x + x}{2}{\Big)} = 2sin{\Big(}\dfrac{\Delta x}{2}{\Big)}\cdot cos\Big(x + \dfrac{ \Delta x }{2}{\Big)}
$$

Тогда

$$
\lim\limits_{\Delta x \to 0}\Delta y = \lim\limits_{\Delta x \to 0}\Bigg(2sin{\Big(}\dfrac{\Delta x}{2}{\Big)}\cdot cos\Big(x + \dfrac{ \Delta x }{2}{\Big)}\Bigg) = 0, 
$$

Т.е. $\lim\limits_{\Delta x \to 0}\Delta y = 0 \overset{\text{по опр.}}\Rightarrow y = sin(x)$ непрерывна на всей числовой прямой.
\vspace*{20pt} 

Рассмотрим $y = cos(x)$:

Найдем приращение функции $$
\Delta y = f(x + \Delta x) - f(x) = cos(x + \Delta x) - cos(x) =
$$ $$
=2sin{\Big(}\dfrac{x + \Delta x + x}{2}{\Big)}\cdot sin\Big(\dfrac{x - \Delta x - x}{2}{\Big)} = -2sin\Big(x + \dfrac{ \Delta x }{2}{\Big)}\cdot sin{\Big(}\dfrac{\Delta x}{2}{\Big)} 
$$

Тогда

$$
\lim\limits_{\Delta x \to 0}\Delta y = \lim\limits_{\Delta x \to 0}\Bigg(-2sin\Big(x + \dfrac{ \Delta x }{2}{\Big)}\cdot sin{\Big(}\dfrac{\Delta x}{2}{\Big)}\Bigg) = 0, 
$$

Т.е. $\lim\limits_{\Delta x \to 0}\Delta y = 0 \overset{\text{по опр.}}\Rightarrow y = cos(x)$ непрерывна на всей числовой прямой.
% определение непр в точке 
\newpage 
\subsection*{20. \textit{Свойства функций, непрерывных на отрезке}}\addcontentsline{toc}{subsection}{20. \textit{Свойства функций, непрерывных на отрезке}}
\begin{enumerate}

\item \textit{Первая теорема Вейерштрасса}

Если функция $y = f(x)$ является непрерывной на $[a, b]$, то она ограничена на этом отрезке.
\item \textit{Вторая теорема Вейерштрасса}

Если функция $y = f(x)$ является непрерывной на $[a, b]$, то она имеет на этом отрезке наибольшее и наименьшее значение. % правильно? 
\item \textit{Первая теорема Больцано-Коши}

Если функция $y = f(x)$ является непрерывной на $[a, b]$ и на концах этого отрезка принимает значения разных знаков, т.е. $f(a)\cdot f(b) < 0$, то существует хотя бы одна точка $c \in (a, b)$, в которой значение функции $f(c) = 0$.
\item \textit{Вторая теорема Больцано-Коши}

Если функция $y = f(x)$ является непрерывной на $[a, b]$ и $f(a)\neq f(b)$, то существует такая точка $c \in (a, b)$, что $f(a) < f(c) < f(b)$.
\item \textit{Теорема о непрерывности обратной функции}

Если функция $y = f(x)$ непрерывна и монотонно возрастает (убывает) на $[a, b]$, то существует и определена на отрезке $[f(a), f(b)]$ обратная функция $x = f^{-1}(y),$ непрерывная и возрастающая (убывающая) на этом отрезке.

\end{enumerate}
\newpage 
\subsection*{21. \textit{Точки разрыва функции и их классификация. Примеры}}\addcontentsline{toc}{subsection}{21. \textit{Точки разрыва функции и их классификация. Примеры}}

Если данная функция $f(x)$ не является непрерывной в точке $x_0$, то $x_0$ называется \textbf{точкой разрыва} функции $f(x)$.
\begin{enumerate}

\item \textbf{Точкой разрыва первого рода} называют такую точку разрыва функции, в которой существуют оба односторонних предела этой функции и они конечны.

Пример: $f(x) = {\dfrac{sin(x)}{x}}$. Рассмотрим односторонние пределы и значение функции в точке $x = 0$:

$$\begin{rcases}
        \exists\lim\limits_{x \to 0+}{\dfrac{sin(x)}{x}} = 1\\
        \exists\lim\limits_{x \to 0-}{\dfrac{sin(x)}{x}} = 1\\
        x = 0 \notin D_f, \text{ т.е. } \nexists f(0) 
    \end{rcases}
    \Rightarrow \text{ точка }x = 0 \text{ - точка разрыва первого рода.}
    $$
\item \textbf{Точкой разрыва второго рода} называют такую точку разрыва функции, в которой хотя бы один из односторонних пределов функции не существует (в частности, равен бесконечности).

Пример: $f(x) = {\dfrac{1}{x}}$. Рассмотрим односторонние пределы и значение функции в точке $x = 0$

$$\begin{rcases}
        \exists\lim\limits_{x \to 0+}{\dfrac{1}{x}} = +\infty\\
        \exists\lim\limits_{x \to 0-}{\dfrac{1}{x}} = -\infty\\
        x = 0 \notin D_f, \text{ т.е. } \nexists f(0) 
    \end{rcases} 
    \Rightarrow \text{ точка }x = 0 \text{ - точка разрыва второго рода.}
    $$
\item Если $x_0$ — точка разрыва функции первого рода и односторонние пределы функции в этой точке равны между собой, но не равны значению функции в этой точке или $f(x)$ не определена в этой точке, то такой разрыв называют \textbf{устранимым}, а точку $x_0$ - \textbf{точкой устранимого разрыва первого рода}.

Пример: $f(x) = {\dfrac{sin(x)}{x}}$:

Из соображений выше, точка $x = 0$ является точкой разрыва первого рода функции $f(x)$. При этом, $\lim\limits_{x \to 0+}{\dfrac{sin(x)}{x}} = \lim\limits_{x \to 0-}{\dfrac{sin(x)}{x}} = 1$, тогда по определению, точка $x = 0$ - точка устранимого разрыва первого рода.

Если доопределить функцию $f(x)$ следующим образом:

$$f(x) = 
    \begin{cases}
        {\dfrac{sin(x)}{x}},\ x \neq 0\\
        1, \hspace*{28pt} x = 0
    \end{cases}
    $$

то функция $f(x)$ по определению будет непрерывной.
\item Если $x_0$ — точка разрыва функции первого рода и односторонние пределы функции в этой точке не равны между собой, то такой разрыв называют \textbf{неустранимым}, а точку $x_0$ - \textbf{точкой неустранимого разрыва первого рода}.

Пример:

$$f(x) = 
    \begin{cases}
        x - 1,\ x \leqslant 0\\
        x + 1, \  x > 0
    \end{cases}
    $$ Рассмотрим односторонние пределы и значение функции в точке $x = 0$:

$$\begin{rcases}
        \exists\lim\limits_{x \rightarrow 0+}f(x) = 1\\
        \exists\lim\limits_{x \rightarrow 0-}f(x) = -1\\
        \exists f(0) = -1
    \end{rcases} 
    \Rightarrow
    $$

$$\Rightarrow \text{ точка }x = 0 \text{ - точка неустранимого разрыва первого рода по определению.}
    $$

\end{enumerate}
\newpage 
\subsection*{22. \textit{Теорема (о необходимом и достаточном условии существования наклонной асимптоты)}}\addcontentsline{toc}{subsection}{22. \textit{Теорема (о необходимом и достаточном условии существования наклонной асимптоты)}}
\begin{Quote2} 
\small\centering 

Используются определения 35, 84, 85 из списка \end{Quote2} 

Прямая $y = kx + b$ является правой (левой) наклонной асимптотой графика функции $y = f(x) \iff$ существуют конечные пределы

$$\lim\limits_{x \to +(-)\infty}{\dfrac{f(x)}{x}} = k,$$ $$\lim\limits_{x \to +(-)\infty}{\big(f(x) - kx\big)} = b$$
\vspace*{20pt} 

\textit{Доказательство}

$(\Rightarrow)$

По условию прямая $y = kx + b$ является правой (левой) наклонной асимптотой графика $y = f(x)$. Тогда по определению

$$
f(x) = kx+b+\underbrace{\alpha(x)}_{\text{б.м.ф.}},\ \text{при } x \rightarrow +(-)\infty
$$

Отсюда $\ {\dfrac{f(x)}{x}} = k + {\dfrac{b}{x}} + {\dfrac{\alpha(x)}{x}} \rightarrow k, \ f(x) - kx = b + \alpha(x), $ при $x \rightarrow +(-)\infty$, т.е. $\lim\limits_{x \to +(-)\infty}{\dfrac{f(x)}{x}} = k,  \lim\limits_{x \to +(-)\infty}{\big(f(x) - kx\big)} = b$

$(\Leftarrow)$

По условию существуют конечные пределы

$$\lim\limits_{x \to +(-)\infty}{\dfrac{f(x)}{x}} = k,$$ $$\lim\limits_{x \to +(-)\infty}{\big(f(x) - kx\big)} = b$$

тогда по теореме о связи функции, ее предела и бесконечно малой

$${f(x) - kx} = b + \underbrace{\alpha(x)}_{\text{б.м.ф.}},\ \text{при } x \rightarrow +(-)\infty$$ $${f(x)} = kx + b + \underbrace{\alpha(x)}_{\text{б.м.ф.}},\ \text{при } x \rightarrow +(-)\infty$$

Таким образом, прямая $y = kx + b$ является правой (левой) наклонной асимптотой графика функции $f(x)$ по определению.

Теорема доказана.
\newpage 
\subsection*{23. \textit{Теорема (необходимое и достаточное условие дифференцируемости функции)}}\addcontentsline{toc}{subsection}{23. \textit{Теорема (необходимое и достаточное условие дифференцируемости функции)}}
\begin{Quote2} 
\small\centering 

Используются определения 61, 64 из списка \end{Quote2} 

Функция $f(x)$ дифференцируема в некоторой точке $x_0$ тогда и только тогда, когда существует конечная производная $f'(x_0)$ в этой точке.
\vspace*{20pt} 

\textit{Доказательство}

$(\Rightarrow)$.

По условию функция $f(x)$ дифференцируема в точке $x_0$. Тогда по определению дифференцируемости:

$$\Delta y  = A \cdot \Delta x + \alpha(\Delta x)\cdot \Delta x,$$

После деления на $\Delta x$ получаем:

$$
\dfrac{\Delta y}{\Delta x}
= A + \alpha(\Delta x), \ \text{где } \lim\limits_{\Delta x \rightarrow 0}\alpha(\Delta x) = 0.
$$

По теореме о связи функции, ее предела и бесконечно малой, имеем

$$
\lim\limits_{\Delta x \rightarrow 0}\dfrac{\Delta y}{\Delta x} = A,
$$

Таким образом производная $f'(x_0)$ существует (и равна A) по определению.

$(\Leftarrow)$

По условию существует $f'(x_0)$. Тогда по определению производной:

$$
\lim\limits_{\Delta x \rightarrow 0}\dfrac{\Delta y}{\Delta x} = f'(x_0)
$$

Тогда по теореме о связи функции, ее предела и бесконечно малой, имеем

$$
\dfrac{\Delta y}{\Delta x} = f'(x_0) + \alpha(\Delta x), \ \text{где } \lim\limits_{\Delta x \rightarrow 0}\alpha(\Delta x) = 0.
$$

После умножения на $\Delta x$ получаем:

$$
\Delta y = f'(x_0)\cdot \Delta x + \alpha(\Delta x)\cdot \Delta x
$$

Таким образом, функция $f(x)$ дифференцируема в точке $x_0$ по определению.

Теорема доказана.
\newpage 
\subsection*{24. \textit{Теорема (о связи дифференцируемости и непрерывности функции)}}\addcontentsline{toc}{subsection}{24. \textit{Теорема (о связи дифференцируемости и непрерывности функции)}}
\begin{Quote2} 
\small\centering 

Используются определения 51, 64 из списка \end{Quote2} 

Если функция $f(x)$ дифференцируема в некоторой точке $x = x_0$, то она непрерывна в этой точке.
\vspace*{20pt} 

\textit{Доказательство}

По условию $y = f(x)$ дифференцируема в точке $x = x_0$, тогда по определению

$$
\Delta y = A \cdot \Delta x + \alpha(\Delta x)\cdot \Delta x, \ \text{где } \lim\limits_{\Delta x \rightarrow 0}\alpha(\Delta x) = 0. 
$$

$$\Downarrow $$

$$
\lim\limits_{\Delta x \rightarrow 0}\Delta y = \lim\limits_{\Delta x \rightarrow 0}\big(A \cdot \Delta x + \alpha(\Delta x)\cdot \Delta x \big) = 0, \overset{\text{по опр.}}{\Rightarrow} y = f(x) \text{ непрерывна в точке } x = x_0
$$

Теорема доказана.
\newpage 
\subsection*{25. \textit{Теорема (о производной произведения двух дифференцируемых функций)}}\addcontentsline{toc}{subsection}{25. \textit{Теорема (о производной произведения двух дифференцируемых функций)}}
\begin{Quote2} 
\small\centering 

Используются определения 61, 64 из списка \end{Quote2} 

Если функции $f(x)$ и $g(x)$ дифференцируемы в точке $x_0$, то функция $f(x)\cdot g(x)$ тоже дифференцируема в этой точке и $\big(f(x)\cdot g(x)\big)' = f'(x)\cdot g(x) + f(x)\cdot g'(x).$
\vspace*{20pt} 

\textit{Доказательство}

По условию функции $f(x)$ и $g(x)$ дифференцируемы в точке $x = x_0, \Rightarrow$ существуют конечные пределы

$$
\lim\limits_{\Delta x \rightarrow 0}\dfrac{\Delta f}{\Delta x} = f'(x_0), \ \lim\limits_{\Delta x \rightarrow 0}\dfrac{\Delta g}{\Delta x} = g'(x_0)
$$

Вычислим $\big(f(x)\cdot g(x)\big)':$ $$
\big(f(x)\cdot g(x)\big)' 
= \lim\limits_{\Delta x \rightarrow 0}{\dfrac{\Delta \big(f(x)\cdot g(x)\big)}{\Delta x}} 
= \lim\limits_{\Delta x \rightarrow 0}{\dfrac{f(x + \Delta x)\cdot g(x + \Delta x) - f(x)\cdot g(x)}{\Delta x}} 
= 
$$

$$
\lim\limits_{\Delta x \rightarrow 0}{\dfrac{f(x + \Delta x)\cdot g(x + \Delta x) - f(x)\cdot g(x + \Delta x) + f(x)\cdot g(x + \Delta x) - f(x)\cdot g(x)}{\Delta x}} 
=
$$

$$
\lim\limits_{\Delta x \rightarrow 0}\bigg(\dfrac{f(x + \Delta x)- f(x)}{\Delta x}\cdot g(x + \Delta x) + f(x)\cdot \dfrac{g(x + \Delta x)- g(x)}{\Delta x}\bigg) 
=
$$

$$
\lim\limits_{\Delta x \rightarrow 0}\bigg(\dfrac{\Delta f}{\Delta x}\cdot g(x + \Delta x) + f(x)\cdot \dfrac{\Delta g}{\Delta x}\bigg)
=\underbrace{\lim\limits_{\Delta x \rightarrow 0}\dfrac{\Delta f}{\Delta x}}_{f'(x)}\cdot \underbrace{\lim\limits_{\Delta x \rightarrow 0}g(x + \Delta x)}_{g(x)} + \underbrace{\lim\limits_{\Delta x \rightarrow 0}f(x)}_{f(x)}\cdot \underbrace{\lim\limits_{\Delta x \rightarrow 0}\dfrac{\Delta g}{\Delta x}}_{g'(x)}=
$$

$$ = f'(x)\cdot g(x) + f(x)\cdot g'(x) $$

Теорема доказана.
\newpage 
\subsection*{26. \textit{Теорема (о производной частного двух дифференцируемых функций)}}\addcontentsline{toc}{subsection}{26. \textit{Теорема (о производной частного двух дифференцируемых функций)}}
\begin{Quote2} 
\small\centering 

Используются определения 61, 64 из списка \end{Quote2} 

Если функции $f(x)$ и $g(x)$ дифференцируемы в точке $x_0$, то функция ${\dfrac{f(x)}{g(x)}}$ тоже дифференцируема в этой точке $\big($при условии $g(x) \neq 0\big)$ и $\bigg({\dfrac{f(x)}{g(x)}}\bigg)' = {\dfrac{f'(x)\cdot g(x) - f(x)\cdot g'(x)}{g^2(x)}}.$
\vspace*{20pt} 

\textit{Доказательство}

По условию функции $f(x)$ и $g(x)$ дифференцируемы в точке $x = x_0, \Rightarrow$ существуют конечные пределы

$$
\lim\limits_{\Delta x \rightarrow 0}\dfrac{\Delta f}{\Delta x} = f'(x_0), \ \lim\limits_{\Delta x \rightarrow 0}\dfrac{\Delta g}{\Delta x} = g'(x_0)
$$

Вычислим $\bigg({\dfrac{f(x)}{g(x)}}\bigg)'$:

$$
\bigg({\dfrac{f(x)}{g(x)}}\bigg)' = \lim\limits_{\Delta x \rightarrow 0} {\dfrac{\Delta\bigg({\dfrac{f(x)}{g(x)}}\bigg)}{\Delta x}}
= \lim\limits_{\Delta x \rightarrow 0}\Bigg( \dfrac{1}{\Delta x}\cdot\bigg({\dfrac{f(x + \Delta x)}{g(x + \Delta x)}}-{\dfrac{f(x)}{g(x)}}\bigg)\Bigg) 
=
$$

$$
=\lim\limits_{\Delta x \rightarrow 0}\Bigg( \dfrac{1}{\Delta x}\cdot\bigg({\dfrac{f(x + \Delta x)\cdot g(x)-g(x + \Delta x)\cdot f(x)}{g(x)\cdot g(x + \Delta x)}}\bigg)\Bigg) =
$$

$$
=\lim\limits_{\Delta x \rightarrow 0}\Bigg( \dfrac{1}{\Delta x}\cdot\bigg({\dfrac{f(x + \Delta x)\cdot g(x) - f(x)\cdot g(x) + f(x)\cdot g(x) -g(x + \Delta x)\cdot f(x)}{g(x)\cdot g(x + \Delta x)}}\bigg)\Bigg) 
=
$$

$$
=\lim\limits_{\Delta x \rightarrow 0}\Bigg( {\dfrac{{\dfrac{f(x + \Delta x) - f(x)}{\Delta x}}\cdot g(x) -  {\dfrac{g(x + \Delta x) - g(x)}{\Delta x}}\cdot f(x)}{g(x)\cdot g(x + \Delta x)}}\Bigg) 
=
$$

$$
=\lim\limits_{\Delta x \rightarrow 0}\Bigg(\bigg({\dfrac{\Delta f}{\Delta x}}\cdot g(x) -  {\dfrac{\Delta g}{\Delta x}}\cdot f(x)\bigg)\cdot\dfrac{1}{g(x)\cdot g(x + \Delta x)}\Bigg) =
$$

$$
=\Bigg(\underbrace{\lim\limits_{\Delta x \rightarrow 0}{\dfrac{\Delta f}{\Delta x}}}_{f'(x)}\cdot \underbrace{\lim\limits_{\Delta x \rightarrow 0}g(x)}_{g(x)} -  \underbrace{\lim\limits_{\Delta x \rightarrow 0}{\dfrac{\Delta g}{\Delta x}}}_{g'(x)}\cdot \underbrace{\lim\limits_{\Delta x \rightarrow 0}f(x)}_{f(x)}\Bigg)\cdot\underbrace{\lim\limits_{\Delta x \rightarrow 0}\dfrac{1}{g(x)\cdot g(x + \Delta x)}}_{g^{-2}(x)}=
$$

$$
= {\dfrac{f'(x)\cdot g(x) - f(x)\cdot g'(x)}{g^2(x)}}
$$

Теорема доказана.
\newpage 
\subsection*{27. \textit{Теорема (о производной сложной функции)}}\addcontentsline{toc}{subsection}{27. \textit{Теорема (о производной сложной функции)}}
\begin{Quote2} 
\small\centering 

Используются определения 51, 61, 64 из списка \end{Quote2} 

Если функция $y=f(u)$ дифференцируема в точке $u_0$ и функция $u = g(x)$ дифференцируема в точке $x_0$, $u_0 = g(x_0)$, то сложная функция $y = f(g(x))$ дифференцируема в точке $x_0$, и $\Big(f\big(g(x)\big)\Big)' = f'_u\cdot g'_x$.
\vspace*{20pt} 

\textit{Доказательство}

По условию
\begin{itemize}

\item функция $u = g(x)$ дифференцируема в точке $x = x_0,$ тогда по определению существует конечный $\lim\limits_{\Delta x\rightarrow 0}{\dfrac{\Delta u}{\Delta x}} = g'(x_0)$
\item функция $y = f(u)$ дифференцируема в точке $u = u_0,$ тогда по определению существует конечный $\lim\limits_{\Delta u\rightarrow 0}{\dfrac{\Delta y}{\Delta u}} = f'(u_0)$
\item функция $u = g(x)$ дифференцируема в точке $x = x_0, \ \Rightarrow$ Функция $u = g(x)$ непрерывна в точке $x = x_0 \ \overset{\text{по опр.}}{\Rightarrow} \lim\limits_{\Delta x \rightarrow 0}\Delta u = 0, \text{ т.е. } \Delta u \rightarrow 0 \text{ при } \Delta x \rightarrow 0$.

\end{itemize}

Таким образом,

$$
y'_x=\lim\limits_{\Delta x \rightarrow 0}\dfrac{\Delta y}{\Delta x}=\lim\limits_{\Delta x \rightarrow 0}\dfrac{\Delta y\cdot \Delta u}{\Delta u \cdot \Delta x} = \lim\limits_{\Delta u \rightarrow 0}\dfrac{\Delta y}{\Delta u}\cdot \lim\limits_{\Delta x \rightarrow 0}\dfrac{\Delta u}{\Delta x} = f'(u_0)\cdot g'(x_0)
$$
\newpage 
\subsection*{28. \textit{Теорема (о производной обратной функции)}}\addcontentsline{toc}{subsection}{28. \textit{Теорема (о производной обратной функции)}}
\begin{Quote2} 
\small\centering 

Используются определения 51, 61 из списка \end{Quote2} 

Если функция $y = f(x)$ строго монотонна и дифференцируема в точке $x = x_0$, то обратная ей функция $x = f^{-1}(y)$ дифференцируема в точке $y = f(x_0)$ и $x'_y = {\dfrac{1}{y'_x}}$.
\vspace*{20pt} 

\textit{Доказательство}

По условию функция $y = f(x)$ дифференцируема в точке $x = x_0, \ \Rightarrow$ существует конечный $\lim\limits_{\Delta x \rightarrow 0}{\dfrac{\Delta y}{\Delta x}} = f'(x_0)$

Функция $y = f(x)$ дифференцируема в точке $x = x_0, \ \Rightarrow$ Функция $y = f(x)$ непрерывна в точке $x = x_0 \ \overset{\text{по опр.}}{\Rightarrow} \lim\limits_{\Delta x \rightarrow 0}\Delta y = 0, \text{ т.е. } \Delta y \rightarrow 0 \text{ при } \Delta x \rightarrow 0$.

Тогда $$
\big(f^{-1}(y)\big)' = x'_y \overset{\text{по опр.}}{=} \lim\limits_{\Delta y \rightarrow 0} {\dfrac{\Delta x}{\Delta y}}=\lim\limits_{\Delta x \rightarrow 0} {\dfrac{\Delta x}{\Delta y}} = \dfrac{1}{\lim\limits_{\Delta x \rightarrow 0} {\dfrac{\Delta y}{\Delta x}}} = \dfrac{1}{y'_x}
$$

Теорема доказана.
\newpage 
\subsection*{29. \textit{Теорема (свойство инвариантности формы записи дифференциала первого порядка)}}\addcontentsline{toc}{subsection}{29. \textit{Теорема (свойство инвариантности формы записи дифференциала первого порядка)}}
\begin{Quote2} 
\small\centering 

Используются определения 65 из списка \end{Quote2} 

Дифференциал функции $y = f(u)$ не зависит от того, является ли $u$ - независимой переменной, или функцией от другой независимой переменной.
\vspace*{20pt} 

\textit{Доказательство}
\begin{enumerate}

\item Пусть $y = f(u),$ где $u$ - независимая переменная. Тогда

$$dy = f'(u)\cdot du
    $$
\item Пусть $y = f(u),$ где $u = g(x)$ - некоторая функция, имеющая производную. Тогда

$$dy = y'_x\cdot du = y'_u\cdot \underbrace{u'_x \cdot dx}_{du} = y'_u\cdot du = f'(u) \cdot du.
    $$

\end{enumerate}

Теорема доказана.
\newpage 
\subsection*{30. \textit{Теорема Ферма}}\addcontentsline{toc}{subsection}{30. \textit{Теорема Ферма}}
\begin{Quote2} 
\small\centering 

Используются определения 61, 64, 75 из списка \end{Quote2} 

Если функция $y = f(x)$ дифференцируема в точке $x_0$, и точка $x_0$ - есть точка локального экстремума, то $f'(x_0) = 0$.
\vspace*{20pt} 

\textit{Доказательство}

Пусть $x_0$ - точка локального максимума функции $y = f(x)$, тогда по определению

$$
\exists U_\delta(x_0) : \ \forall x \in U_\delta(x_0) \Rightarrow f(x) \leqslant f(x_0)
$$

По условию $y = f(x)$ дифференцируема в точке $x = x_0, \ \Rightarrow$ в точке $x = x_0$ существует конечная производная $\lim\limits_{\Delta x \rightarrow 0}{\dfrac{\Delta y}{\Delta x}} = f'(x_0)$

Тогда

$$
f(x_0 + \Delta x) \leqslant f(x_0), \ x_0 + \Delta x \in U_\delta(x_0)
$$

$$
f(x_0 + \Delta x) - f(x_0) \leqslant 0
$$

$$
\text{Если } \Delta x > 0, \text{ то } {\dfrac{f(x_0+\Delta x) - f(x_0)}{\Delta x}} \leqslant 0
$$

$$
\text{Если } \Delta x < 0, \text{ то }{\dfrac{f(x_0+\Delta x) - f(x_0)}{\Delta x}} \geqslant 0
$$

По теореме о переходе к пределу в неравенстве

$$
\text{Если } \Delta x > 0, \text{ то } \lim\limits_{\Delta x \rightarrow 0}{\dfrac{f(x_0 + \Delta x) - f(x_0)}{\Delta x}} = f'(x_0) \leqslant 0
$$

$$
\text{Если } \Delta x < 0, \text{ то }\lim\limits_{\Delta x \rightarrow 0}{\dfrac{f(x_0 + \Delta x) - f(x_0)}{\Delta x}} = f'(x_0) \geqslant 0
$$

Таким образом, $f'(x_0) = 0$.

Теорема доказана.
\newpage 
\subsection*{31. \textit{Теорема Ролля}}\addcontentsline{toc}{subsection}{31. \textit{Теорема Ролля}}
\begin{Quote2} 
\small\centering 

Используются определения 55, 64, 74, 75 из списка \end{Quote2} 

Пусть функция $y = f(x):$
\begin{enumerate}

\item Непрерывна на отрезке $[a, b]$
\item Дифференцируема на интервале $(a, b)$
\item $f(a) = f(b)$

\end{enumerate}

Тогда на интервале $(a, b)$ существует по крайней мере одна точка $x_0,$ в которой $f'(x_0) = 0$.
\vspace*{20pt} 

\textit{Доказательство}

По условию $y = f(x)$ непрерывна на $[a, b], \Rightarrow $ по (второй) теореме Вейерштрасса функция $y = f(x)$ на отрезке $[a,b]$ достигает своего наибольшего и наименьшего значения, обозначим

$$
M = \underset{[a, b]}{max}\big(f(x)\big)
$$

$$
m = \underset{[a, b]}{min}\big(f(x)\big)
$$

тогда $\forall x \in [a, b] : \ m \leqslant f(x) \leqslant M$

Если $m = M$, то $\forall x \in [a, b] \ \ m = M = f(x) = const \Rightarrow \forall x \in [a, b] \ \ f'(x) = 0$

Если $m \neq M$, то
\begin{itemize}

\item $f(a) = f(b) = m$. Тогда функция $y = f(x)$ достигает своего наибольшего значения внутри $[a, b]$, т.е. $a < x_0 < b$

Таким образом, точка $x_0$ - точка локального максимума. Также по условию, $f(x)$ дифференцируема на интервале $(a, b) \Rightarrow$ $f(x)$ дифференцируема в точке $x_0$. В итоге, по теореме Ферма, $f'(x_0)=0$
\item $f(a) = f(b) = M$. Тогда функция $y = f(x)$ достигает своего наименьшего значения внутри $[a, b]$, т.е. $a < x_0 < b$

Таким образом, точка $x_0$ - точка локального минимума. Также по условию, $f(x)$ дифференцируема на интервале $(a, b) \Rightarrow$ $f(x)$ дифференцируема в точке $x_0$. В итоге, по теореме Ферма, $f'(x_0)=0$
\item $y = f(x)$ достигает своего минимального и максимального значения внутри $[a, b]$ в точках $x_0$ и $ x_1$.

\end{itemize}

Точки $x_0$ и $ x_1$ - точки экстремума. Также по условию, $f(x)$ дифференцируема на интервале $(a, b) \Rightarrow$ $f(x)$ дифференцируема в точках $x_0$ и $x_1$. В итоге, по теореме Ферма, $f'(x_0)=0$, $f'(x_1)=0$.

Теорема доказана.
\newpage 
\subsection*{32. \textit{Теорема Лагранжа}}\addcontentsline{toc}{subsection}{32. \textit{Теорема Лагранжа}}
\begin{Quote2} 
\small\centering 

Используются определения 55, 64 из списка \end{Quote2} 

Пусть функция $f(x)$:
\begin{enumerate}

\item Непрерывна на отрезке $[a, b]$
\item Дифференцируема на интервале $(a, b)$

\end{enumerate}

Тогда существует хотя бы одна точка $c \in (a, b),$ такая, что $f(b) - f(a) = f'(c)\cdot(b - a)$
\vspace*{20pt} 

\textit{Доказательство}

Рассмотрим вспомогательную функцию $$F(x) = f(x) - \dfrac{f(b) - f(a)}{b - a}\cdot (x - a)$$
\begin{enumerate}

\item Эта функция непрерывна на отрезке $[a, b]$ и дифференцируема на интервале $(a, b)$, поскольку этими свойствами обладает $f(x)$.
\item $$\begin{rcases}
    F(a) =f(a) - {\dfrac{f(b) - f(a)}{b - a}}\cdot (a - a) = f(a)\\
    F(b) = f(b) - {\dfrac{f(b) - f(a)}{b - a}}\cdot (b - a) = f(a)
    \end{rcases} \Rightarrow F(a) = F(b)
    $$

\end{enumerate}

Таким образом, для $F(x)$ выполнены все условия теоремы Ролля $\Rightarrow$ существует точка $c \in (a, b)$, для которой $F'(c) = f'(c) - {\dfrac{f(b) - f(a)}{b - a}}= 0$

$$f'(c) = {\dfrac{f(b) - f(a)}{b - a}}$$ $$f'(c)\cdot (b - a) = f(b) - f(a)$$

Теорема доказана.
\newpage 
\subsection*{33. \textit{Теорема Коши}}\addcontentsline{toc}{subsection}{33. \textit{Теорема Коши}}
\begin{Quote2} 
\small\centering 

Используются определения 55, 64 из списка \end{Quote2} 

Пусть функции $f(x)$ и $g(x)$:
\begin{enumerate}

\item Непрерывны на отрезке $[a, b]$
\item Дифференцируемы на интервале $(a, b)$
\item $g'(x) \neq 0, \ \forall x \in (a, b)$

\end{enumerate}

Тогда существует хотя бы одна точка $c \in (a, b),$ такая, что

$$
\dfrac{f(b) - f(a)}{g(b)-g(a)} = \dfrac{f'(c)}{g'(c)}
$$
\vspace*{20pt} 

\textit{Доказательство}

Сначала заметим, что $g(b) - g(a) \neq 0$, т.к. если бы $g(b) - g(a) = 0$, то $g(a) = g(b)$ и функция $g(x)$, в результате, удовлетворяла бы условию теоремы Ролля, согласно которой $\exists c \in (a, b)$, такая, что $g'(c) = 0$, что противоречит 3 условию теоремы.

Введем вспомогательную функцию

$$
F(x) = f(x) - \dfrac{f(b) - f(a)}{g(b)-g(a)}\cdot \big(g(x) - g(a)\big)
$$
\begin{enumerate}

\item Эта функция непрерывна на отрезке $[a, b]$ и дифференцируема на интервале $(a, b)$, поскольку этими свойствами обладают $f(x)$ и $g(x)$.
\item $$\begin{rcases}
        F(a) = f(a) - {\dfrac{f(b) - f(a)}{g(b)-g(a)}}\cdot \big(g(a) - g(a)\big) = f(a)\\
        F(b) = f(b) - {\dfrac{f(b) - f(a)}{g(b)-g(a)}}\cdot \big(g(b) - g(a)\big) = f(a)
    \end{rcases}
    \Rightarrow F(a) = F(b)
    $$

\end{enumerate}

Таким образом, для $F(x)$ выполнены все условия теоремы Ролля $\Rightarrow$ существует точка $c \in (a, b)$, для которой $F'(c) = f'(c) - {\dfrac{f(b) - f(a)}{g(b)-g(a)}}\cdot g'(c) = 0$

$$f'(c) - {\dfrac{f(b) - f(a)}{g(b)-g(a)}}\cdot g'(c) = 0$$ $$f'(c) = {\dfrac{f(b) - f(a)}{g(b)-g(a)}}\cdot g'(c)$$ $$\dfrac{f'(c)}{g'(c)}=\dfrac{f(b) - f(a)}{g(b)-g(a)}$$

Теорема доказана.
% правильно? 
\newpage 
\subsection*{34. \textit{Теорема Лопиталя – Бернулли для предела отношения двух бесконечно малых функций}}\addcontentsline{toc}{subsection}{34. \textit{Теорема Лопиталя – Бернулли для предела отношения двух бесконечно малых функций}}
\begin{Quote2} 
\small\centering 

Используются определения 35, 36, 64 из списка \end{Quote2} 

Пусть $f(x)$ и $g(x)$:
\begin{enumerate}

\item Являются бесконечно малыми или бесконечно большими функциями при $x \rightarrow x_0$
\item Дифференцируемы в $\overset{\circ}U(x_0)$
\item $g'(x) \neq 0$ в $\overset{\circ}U(x_0)$
\item $\exists \lim\limits_{x \rightarrow x_0}{\dfrac{f'(x)}{g'(x)}}$

\end{enumerate}

Тогда существует предел $\lim\limits_{x \rightarrow x_0}{\dfrac{f(x)}{g(x)}} = \lim\limits_{x \rightarrow x_0}{\dfrac{f'(x)}{g'(x)}}$
\vspace*{20pt} 

\textit{Доказательство (для б.м.ф)}

По условию $f(x)$ и $g(x)$ являются бесконечно малыми при $x \rightarrow x_0$, тогда по определению

$$
\lim\limits_{x \rightarrow x_0}f(x) = 0 \iff \lim\limits_{x \rightarrow x_0+}f(x) = \lim\limits_{x \rightarrow x_0-}f(x) = 0
$$

$$
\lim\limits_{x \rightarrow x_0}g(x) = 0 \iff \lim\limits_{x \rightarrow x_0+}g(x) = \lim\limits_{x \rightarrow x_0-}g(x)=0
$$

По условию $\exists \lim\limits_{x \rightarrow x_0}{\dfrac{f'(x)}{g'(x)}} \iff \lim\limits_{x \rightarrow x_0+}{\dfrac{f'(x)}{g'(x)}} = \lim\limits_{x \rightarrow x_0-}{\dfrac{f'(x)}{g'(x)}} = \lim\limits_{x \rightarrow x_0}{\dfrac{f'(x)}{g'(x)}}$

Рассмотрим пределы при $x \rightarrow x_0+$ (для $x \rightarrow x_0- \ $ доказывается аналогично)

Доопределим функции $f(x)$ и $g(x)$ в точке $x = x_0,$ полагая $f(x_0) = g(x_0) = 0$. Тогда $f(x)$ и $g(x)$ определены и непрерывны в $U^+(x_0)$. Рассмотрим отрезок $[x_0, x],$ где $x > x_0$.

По условию $g'(x) \neq 0 $ в $(x_0, x)$.

Тогда функции $f(x)$ и $g(x)$ удовлетворяют условию теоремы Коши, а значит, $\exists c \in (x_0, x)$, такая, что ${\dfrac{f(x) - f(x_0)}{g(x) - g(x_0)}} = {\dfrac{f'(c)}{g'(c)}}$
\begin{enumerate}

\item Так как $f(x_0) = g(x_0) = 0,$ то ${\dfrac{f(x)}{g(x)}} = {\dfrac{f'(c)}{g'(c)}}$
\item Так как $x_0 < c < x$, то, если $x \rightarrow x_0 +$, то и $c \rightarrow x_0+$

\end{enumerate}

Таким образом,

$$
\lim\limits_{x \rightarrow x_0+}{\dfrac{f(x)}{g(x)}} = \lim\limits_{x \rightarrow x_0+}{\dfrac{f'(c)}{g'(c)}} = \lim\limits_{c \rightarrow x_0+}{\dfrac{f'(c)}{g'(c)}} =\lim\limits_{x \rightarrow x_0+}{\dfrac{f'(x)}{g'(x)}}, \text{ т.е. } \lim\limits_{x \rightarrow x_0+}{\dfrac{f(x)}{g(x)}}=\lim\limits_{x \rightarrow x_0+}{\dfrac{f'(x)}{g'(x)}}
$$

Аналогично $\lim\limits_{x \rightarrow x_0-}{\dfrac{f(x)}{g(x)}}=\lim\limits_{x \rightarrow x_0-}{\dfrac{f'(x)}{g'(x)}}$

Таким образом, $\exists\lim\limits_{x \rightarrow x_0}{\dfrac{f(x)}{g(x)}}=\lim\limits_{x \rightarrow x_0}{\dfrac{f'(x)}{g'(x)}}$.

Теорема доказана.
% спросить про ( 1 n > 0 ) -->
\newpage 
\subsection*{35. \textit{Сравнение роста показательной, степенной и логарифмической функций на бесконечности.}}\addcontentsline{toc}{subsection}{35. \textit{Сравнение роста показательной, степенной и логарифмической функций на бесконечности.}}
\begin{Quote2} 
\small\centering 

Используются определения 36 из списка \end{Quote2} 
\begin{enumerate}

\item Сравним рост показательной функции $y = a^x$ и степенной функции $y = x^n$ ($a > 1$, $n > 0$):

Так как при $x \rightarrow +\infty$ функции $y = a^x$ и $y = x^n$ являются бесконечно большими, воспользуемся правилом Лопиталя-Бернулли $n$ раз:

$$\lim\limits_{x \rightarrow +\infty}{\dfrac{a^x}{x^n}} = \lim\limits_{x \rightarrow +\infty}{\dfrac{ln(a)\cdot a^x}{n\cdot x^{n - 1}}} =...=\lim\limits_{x \rightarrow +\infty}{\dfrac{ln^n(a)\cdot a^x}{n!}} = +\infty
    $$

Таким образом, показательная функция $y = a^x$ растет быстрее степенной функции $y = x^n$.
\item Сравним рост логарифмической функции $y = log_a(x)$ и степенной функции $y = x^n (a > 1, n > 0)$:

Так как при $x \rightarrow +\infty$ функции $y  = log_a(x)$ и $y = x^n$ являются бесконечно большими, воспользуемся правилом Лопиталя-Бернулли:

$$\lim\limits_{x \rightarrow +\infty}{\dfrac{ log_a(x)}{x^n}} 
    = \lim\limits_{x \rightarrow +\infty}{\ \dfrac{\dfrac{1}{x\cdot ln(a)}}{n\cdot x^{n - 1}}} 
    =\dfrac{1}{n\cdot ln(a)}\cdot \lim\limits_{x \rightarrow +\infty}{\ \dfrac{1}{x^{n}}}
    = 0
    $$

Таким образом, степенная функция $y = a^x$ растет быстрее логарифмической функции $y = x^n$.

\textit{Вывод}:

Показательная функция растет быстрее степенной, а степенная - быстрее логарифмической.

\end{enumerate}
\newpage 
\subsection*{36. \textit{Вывод формулы Тейлора с остаточным членом в форме Лагранжа}}\addcontentsline{toc}{subsection}{36. \textit{Вывод формулы Тейлора с остаточным членом в форме Лагранжа}}

Пусть функция $f(x)$ определена в окрестности $U(x_0)$ точки $x_0$ и имеет в этой окрестности производные всех порядков до $(n + 1)$-го включительно. Тогда для любого $x \in U(x_0)$ справедливо равенство:
\begin{mainQuote}

$$ f(x) = \sum\limits_{k=0}^n\dfrac{f^{(k)}(x_0)}{k!}\cdot (x - x_0)^k + \underbrace{\dfrac{f^{(n+1)}\big(x_0 + \theta(x - x_0)\big)}{(n+1)!}\cdot (x - x_0)^{n + 1}}_{ R_{n+1}},  $$

$$ \text{ где } \theta \in (0, 1), \ R_{n+1} - \text{ остаточный член в форме Лагранжа}. $$
\end{mainQuote}

Пусть $x \in U(x_0)$, и пусть для определенности $x > x_0$. Рассмотрим на отрезке $[x0, x]$ две функции

$$
\varphi(t) = f(x) - \sum\limits_{k=0}^n\dfrac{f^{(k)}(t)}{k!}\cdot (x - t)^k
$$ $$
\psi(t) = (x - t)^{n+1}
$$

Для этих функций имеем

$$
\varphi(x) = f(x) - \sum\limits_{k=0}^n\dfrac{f^{(k)}(x)}{k!}\cdot (x - x)^k = f(x) - f(x) = 0, 
$$ $$
\varphi(x_0) = f(x) - \sum\limits_{k=0}^n\dfrac{f^{(k)}(x_0)}{k!}\cdot (x -  x_0)^k,
$$ $$\psi(x) = (x - x)^{n+1} = 0,$$ $$\psi(x_0) = (x - x_0)^{n+1}.$$

Вычислим производные

$$
\varphi'(t) = \bigg(f(x) - f(t) - \sum\limits_{k=1}^n\dfrac{f^{(k)}(t)}{k!}\cdot (x - t)^k\bigg)' = $$ $$=-f'(t) - \sum\limits_{k=1}^n \dfrac{1}{k!}\cdot\bigg(f^{(k+1)}(t) \cdot (x - t)^k - k\cdot f^{(k)}(t)\cdot(x - t)^{k-1}\bigg) =
$$ $$
= -f'(t) - \sum\limits_{k=1}^n \dfrac{f^{(k+1)}(t) }{k!}\cdot (x - t)^k + \sum\limits_{k=1}^n\dfrac{f^{(k)}(t)}{(k - 1)!}\cdot(x - t)^{k-1}  
$$

В последней сумме введем новый индекс суммирования $l = k - 1$. Тогда

$$
\sum\limits_{k=1}^n\dfrac{f^{(k)}(t)}{(k - 1)!}\cdot(x - t)^{k-1} = \sum\limits_{l=0}^{n-1}\dfrac{f^{(l+1)}(t)}{l!}\cdot(x - t)^{l}= f'(t)+ \sum\limits_{l=1}^{n-1}\dfrac{f^{(l+1)}(t)}{l!}\cdot(x - t)^{l}
$$

Следовательно

$$
\varphi'(x) = -f'(t) -\dfrac{f^{(n+1)}(t)}{n!}
\cdot (x - t)^n - \sum\limits_{k=1}^{n-1} \dfrac{f^{(k+1)}(t) }{k!}\cdot (x - t)^k + f'(t) +
$$ $$
+ \sum\limits_{l=1}^{n-1} \dfrac{f^{(l+1)}(t) }{l!}\cdot (x - t)^l = -\dfrac{f^{(n+1)}(t)}{n!} \cdot (x - t)^n,
$$

Т.е. $$
\varphi'(x) =-\dfrac{f^{(n+1)}(t)}{n!} \cdot (x - t)^n
$$

Далее, $\ \psi'(t) = -(n + 1)\cdot (x - t)^n$ , и непосредственно видно, что производная $\psi'(t)$ на интервале $(x_0, x)$ отлична от нуля. К паре функций $\psi(t)$ и $\varphi(t)$ на отрезке $[x0, x]$ применим теорему Коши. Имеем

$$
\dfrac{\varphi(x_0) - \varphi(x)}{\psi(x_0) - \psi(x)} =\dfrac{\varphi'\big(x_0 + \theta(x - x_0)\big)}{\psi'\big(x_0+\theta(x - x_0)\big)} , \text{ где } \theta \in (0, 1).
$$

Таким образом, $0 < \theta<1 \ \iff \ 0 < \theta (x - x_0) < x - x_0 \ \iff \ x_0 < x_0 + \theta (x - x_0) < x, \Rightarrow c = x_0 + \theta (x - x_0) \in (x_0, x)$

Учитывая результаты проведенных вычислений, получаем отсюда:

$$
\dfrac{f(x) - \sum\limits_{k=0}^n{\dfrac{f^{(k)}(x_0)}{k!}}\cdot (x - x_0)^k}{(x - x_0)^{n + 1}} = - \dfrac{f^{(n + 1)}\big(x_0 + \theta(x - x_0)\big)}{n!}\cdot \big(x - x_0 - \theta(x - x_0)\big)^n\times
$$ $$
\times \dfrac{1}{-(n + 1)\big(x - x_0 - \theta(x - x_0)\big)^n} = \dfrac{f^{(n+1)}\big(x_0 + \theta(x - x_0)\big)}{(n+1)!},
$$

Т.е.

$$
f(x) - \sum\limits_{k=0}^n{\dfrac{f^{(k)}(x_0)}{k!}}\cdot (x - x_0)^k = {\dfrac{f^{(n+1)}\big(x_0 + \theta(x - x_0)\big)}{(n+1)!}}\cdot (x - x_0)^{n + 1}
$$

Из последнего равенства следует утверждение теоремы при $x > x_0$. При $x < x_0$ рассуждения аналогичны; если $x = x_0$, то утверждение теоремы очевидно.

Теорема доказана.
\newpage 
\subsection*{37. \textit{Вывод формулы Тейлора с остаточным членом в форме Пеано}}\addcontentsline{toc}{subsection}{37. \textit{Вывод формулы Тейлора с остаточным членом в форме Пеано}}

Пусть функция $f(x)$ определена в окрестности точки $x_0$ и имеет в этой точке производные всех порядков до $n$-го включительно. Тогда справедливо равенство
\begin{mainQuote}

$$ f(x) = \sum\limits_{k=0}^n\dfrac{f^{(k)}(x_0)}{k!}\cdot (x - x_0)^k + \underbrace{o\big((x - x_0)^n\big)}_{R_{n+1}}, \ x \rightarrow x_0, \text{где} $$

$$ R_{n+1} - \text{остаточный член в форме Пеано}$$
\end{mainQuote}

Равенство, которое требуется доказать, означает, что

$$
\lim\limits_{x \rightarrow x_0}{\dfrac{f(x) -\sum\limits_{k=0}^n{\dfrac{f^{(k)}(x_0)}{k!}}\cdot (x - x_0)^k}{(x - x_0)^n}} = 0
$$

Мы имеем здесь дело с неопределенностью ${\{\dfrac{0}{0}\}}$. Чтобы раскрыть её, применим n - 1 раз правило Лопиталя-Бернулли

$$
\lim\limits_{x \rightarrow x_0}{\dfrac{f(x) -\sum\limits_{k=0}^n{\dfrac{f^{(k)}(x_0)}{k!}}\cdot (x - x_0)^k}{(x - x_0)^n}} 
= \lim\limits_{x \rightarrow x_0}\dfrac{f'(x) -\sum\limits_{k=1}^n{\dfrac{f^{(k)}(x_0)}{(k-1)!}}\cdot (x - x_0)^{k-1}}{n(x - x_0)^{n-1}} 
=
$$ $$
= \lim\limits_{x \rightarrow x_0}{\dfrac{f''(x) -\sum\limits_{k=2}^n{\dfrac{f^{(k)}(x_0)}{(k-2)!}}\cdot (x - x_0)^{k-2}}{n(n-1)(x - x_0)^{n-2}}} 
= ...=
$$ $$
=\lim\limits_{x \rightarrow x_0}{\dfrac{f^{(n-1)}(x)-f^{(n-1)}(x_0) - f^{(n)}(x_0)(x - x_0)}{n!(x - x_0)}} 
=
$$ $$
= \dfrac{1}{n!}\lim\limits_{x \rightarrow x_0}\bigg({\dfrac{f^{(n-1)}(x)-f^{(n-1)}(x_0)}{(x - x_0)}}- f^{(n)}(x_0)\bigg)=0,
$$

т.к. $\lim\limits_{x \rightarrow x_0}\bigg({\dfrac{f^{(n-1)}(x)-f^{(n-1)}(x_0)}{(x - x_0)}}\bigg)=f^{(n)}(x_0)$.

Теорема доказана.
\newpage 
\subsection*{38. \textit{Формула Маклорена для функции $y = e^x$ с остаточным членом в форме Лагранжа}}\addcontentsline{toc}{subsection}{38. \textit{Формула Маклорена для функции $y = e^x$ с остаточным членом в форме Лагранжа}}

Найдем производные функции $y = e^x$ до $n$-го порядка:

$$f'(x) = f''(x) = f'''(x) =...=f^{(n)}(x) = e^x$$ $$f'(0) = f''(0) = f'''(0) =...=f^{(n)}(0) = 1$$

Таким образом, получаем
\begin{mainQuote}

$$ e^x = 1 + x + \dfrac{x^2}{2!}+ \dfrac{x^3}{3!}+ \dfrac{x^4}{4!} + ... + \dfrac{x^n}{n!} + \underbrace{\dfrac{x^{n + 1}}{(n+1)!}\cdot e^{\theta x}}_{\text{остаточный\ член}}, \ \theta \in (0, 1) $$
\end{mainQuote}
\newpage 
\subsection*{39. \textit{Формула Маклорена для функции $y = sin(x)$ с остаточным членом в форме Лагранжа}}\addcontentsline{toc}{subsection}{39. \textit{Формула Маклорена для функции $y = sin(x)$ с остаточным членом в форме Лагранжа}}

Найдем производные функции $y = sin(x)$ до $n$-го порядка:

$$f(x) = sin(x), \ f(0) = 0$$ $$f'(x) = cos(x) = sin(x + \dfrac{\pi}{2}), \ f'(0)=1$$ $$f''(x) = -sin(x) = sin(x + 2\cdot \dfrac{\pi}{2}), \ f'(0)=0$$ $$f'''(x) = -cos(x) = sin(x + 3\cdot \dfrac{\pi}{2}), \ f'(0) = -1$$ $$\dots$$ $$f^{(2n + 1)}(x) = (-1)^{n}\cdot cos(x), \ f^{(2n + 1)}(0) = (-1)^{n}$$ $$f^{(2n + 2)}(x) = (-1)^{n+1}\cdot sin(x), \ f^{(2n + 2)}(0) = 0$$

Таким образом, получаем
\begin{mainQuote}

$$ sin(x) = x - \dfrac{x^3}{3!} + \dfrac{x^5}{5!}- \dfrac{x^7}{7!} + \dfrac{x^9}{9!}+...+(-1)^{n}\dfrac{x^{2n+1}}{(2n + 1)!} + $$

$$ + \underbrace{\dfrac{x^{2n+2}}{(2n + 2)!}\cdot sin\bigg(\theta x + (2n + 2)\cdot\dfrac{\pi}{2}\bigg)}_{\text{остаточный \ член}}, \ \theta \in (0, 1) $$
\end{mainQuote}
\newpage 
\subsection*{40. \textit{Формула Маклорена для функции $y = cos(x)$ с остаточным членом в форме Лагранжа}}\addcontentsline{toc}{subsection}{40. \textit{Формула Маклорена для функции $y = cos(x)$ с остаточным членом в форме Лагранжа}}

Найдем производные функции $y = cos(x)$ до $n$-го порядка:

$$f(x) = cos(x), \ f(0) = 1$$ $$f'(x) = -sin(x) = cos(x + \dfrac{\pi}{2}), \ f'(0)=0$$ $$f''(x) = -cos(x) = cos(x + 2\cdot \dfrac{\pi}{2}), \ f'(0)=-1$$ $$f'''(x) = sin(x) = cos(x + 3\cdot \dfrac{\pi}{2}), \ f'(0) = 0$$ $$\dots$$ $$f^{(2n)}(x) = (-1)^{n}\cdot cos(x), \ f^{(2n)}(0) = (-1)^{n}$$ $$f^{(2n + 1)}(x) = (-1)^{n+1}\cdot sin(x), \ f^{(2n + 1)}(0) = 0$$

Таким образом, получаем
\begin{mainQuote}

$$ cos(x) = 1 - \frac{x^2}{2!} + \frac{x^4}{4!}- \frac{x^6}{6!} + \frac{x^8}{8!}+...+(-1)^{n}\frac{x^{2n}}{(2n)!} + $$

$$ + \underbrace{\frac{x^{2n+1}}{(2n + 1)!}\cdot cos\bigg(\theta x + (2n + 1)\cdot\frac{\pi}{2}\bigg)}_{\text{остаточный \ член}},\  \theta \in (0, 1) $$
\end{mainQuote}
\newpage 
\subsection*{41. \textit{Формула Маклорена для функции $y = ln(1 + x)$ с остаточным членом в форме Лагранжа}}\addcontentsline{toc}{subsection}{41. \textit{Формула Маклорена для функции $y = ln(1 + x)$ с остаточным членом в форме Лагранжа}}

Найдем производные функции $y = ln(1+x)$ до $n$-го порядка:

$$f(0) = 0$$ $$f'(x) = \dfrac{1}{1+x}, \ f'(0) = 1$$ $$f''(x) = -\dfrac{1}{(1+x)^2}, \ f''(0) = -1 = -1!$$ $$f'''(x) = \dfrac{2}{(1+x)^3}, \ f'''(0) = 2 = 2!$$ $$f^{IV}(x) = \dfrac{-3\cdot 2}{(1+x)^4}, \ f^{IV}(0) = -3\cdot 2 = -3!$$ $$\dots$$ $$f^{(n)}(x) = \dfrac{(n - 1)!}{(1 + x)^n}\cdot (-1)^{n - 1}, \ f^{(n)}(0)=(-1)^{n-1}(n-1)! $$ $$f^{(n+1)}(x) = \dfrac{n!}{(1 + x)^{n+1}}\cdot (-1)^{n}, \ f^{(n+1)}(0)=(-1)^{n}\cdot n!$$

Таким образом, получаем
\begin{mainQuote}

$$ ln(1 + x) = x - \dfrac{x^2}{2}+ \dfrac{x^3}{3}- \dfrac{x^4}{4}+ \dfrac{x^5}{5} + ... + (-1)^{n - 1}\cdot \dfrac{x^n}{n} + \underbrace{(-1)^n\dfrac{x^{n+1}}{(n+1)(1 + \theta x)^{n+1}}}_{\text{остаточный \ член}}, $$

$$ \theta \in (0, 1)$$
\end{mainQuote}
\newpage 
\subsection*{42. \textit{Формула Маклорена для функции $y = (1 + x)^a$ с остаточным членом в форме Лагранжа}}\addcontentsline{toc}{subsection}{42. \textit{Формула Маклорена для функции $y = (1 + x)^a$ с остаточным членом в форме Лагранжа}}

Найдем производные функции $y = (1 + x)^a$ до $n$-го порядка: $$f(x) = (1+x)^a , \ f(0) = 1$$ $$f'(x)=a \cdot (1+x)^{a - 1}, \ f'(0) = a$$ $$f''(x)=a \cdot (a - 1) \cdot (1+x)^{a - 2}, \ f''(0) = a\cdot (a - 1)$$ $$...$$ $$f^{(n)}(x) = a\cdot(a - 1)\cdot(a-2)\cdot...\cdot(a - (n - 1))\cdot (1+x)^{a-n},$$ $$ f^{(n)}(0) = a\cdot(a - 1)\cdot(a-2)\cdot...\cdot(a - (n - 1))$$

Таким образом, получаем
\begin{mainQuote}

$$ (1+x)^a = 1 + \dfrac{a}{1!}x + \dfrac{a(a-1)}{2!}x^2+ \dfrac{a(a-1)(a-2)}{3!}x^3+ \dfrac{a(a-1)(a-2)(a-3)}{4!}x^4+... $$

$$ +\dfrac{a(a-1)(a-2)...(a-(n-1))}{n!}x^n + \underbrace{\dfrac{a(a-1)...(a -n)}{(n+1)!}\cdot (1+\theta x)^{a-(n+1)}\cdot x^{n + 1}}_{\text{остаточный \ член}},$$

$$ \theta \in (0, 1)$$
\end{mainQuote}
\newpage 
\subsection*{43. \textit{Необходимое и достаточное условие неубывания дифференцируемой функции}}\addcontentsline{toc}{subsection}{43. \textit{Необходимое и достаточное условие неубывания дифференцируемой функции}}
\begin{Quote2} 
\small\centering 

Используются определения 64, 71 из списка \end{Quote2} 

Пусть функция $f(x)$ дифференцируема на интервале $(a, b)$. Для того, чтобы эта функция была неубывающей на интервале $(a, b)$, необходимо и достаточно, чтобы производная $f'(x)$ была неотрицательна $\forall x \in (a, b).$
\vspace*{20pt} 

\textit{Доказательство}

$(\Rightarrow)$

По условию $f(x)$ не убывает на интервале $(a, b)$. Тогда в точке $x \in (a, b)$, в которой функция $f(x)$ дифференцируема, имеем
\begin{itemize}

\item $\Delta x > 0\Rightarrow
f(x + \Delta x) \geqslant  f(x)
\Rightarrow f'(x) = f'_+(x) = \lim\limits_{\Delta x \to 0+ }{\dfrac{f(x + \Delta x) - f(x)}{\Delta x}} \geqslant 0$
\item $\Delta x < 0\Rightarrow
f(x)  \geqslant f(x + \Delta x)
\Rightarrow f'(x) = f'_-(x) = \lim\limits_{\Delta x \to 0-}{\dfrac{f(x + \Delta x) - f(x)}{\Delta x}} \geqslant 0$

\end{itemize}

Таким образом, $\forall x \in (a, b) \Rightarrow f'(x) \geqslant 0$.

$(\Leftarrow)$

По условию во всех точках интервала $(a,b)$, в которых $f(x)$ дифференцируема, выполняется неравенство $f'(x) \geqslant 0$. Пусть $x_1$ и $x_2,\ a < x_1 < x_2 < b,$ — произвольные точки этого промежутка. Тогда функция $y = f(x)$ удовлетворяет теореме Лагранжа, $\Rightarrow  \exists c \in (x_1, x_2)$, такая, что $f(x_2) - f(x_1) = f'(c)(x_2 - x_1)$

Т.к. $\forall x \in (a, b) \Rightarrow f'(x) \geqslant 0$, и $x_2 > x_1$, то $f(x_2) - f(x_1) \geqslant 0 \iff f(x_2) \geqslant  f(x_1) $

А значит, функция $y = f(x)$ - неубывающая на $(a, b)$ по определению.

Теорема доказана.
\newpage 
\subsection*{44. \textit{Необходимое и достаточное условие невозрастания дифференцируемой функции}}\addcontentsline{toc}{subsection}{44. \textit{Необходимое и достаточное условие невозрастания дифференцируемой функции}}
\begin{Quote2} 
\small\centering 

Используются определения 64, 69 из списка \end{Quote2} 

Пусть функция $f(x)$ дифференцируема на интервале $(a, b)$. Для того, чтобы эта функция была невозрастающей на интервале $(a, b)$, необходимо и достаточно, чтобы производная $f'(x) \leqslant 0 \ \ \forall x \in (a, b).$
\vspace*{20pt} 

\textit{Доказательство}

$(\Rightarrow)$

По условию $f(x)$ не возрастает на интервале $(a, b)$. Тогда в точке $x \in (a, b)$, в которой функция $f(x)$ дифференцируема, имеем
\begin{itemize}

\item $\Delta x > 0\Rightarrow f(x + \Delta x) \leqslant f(x) \Rightarrow f'(x) = f'_+(x) = \lim\limits_{\Delta x \to 0+ }{\dfrac{f(x + \Delta x) - f(x)}{\Delta x}} \leqslant 0$
\item $\Delta x < 0\Rightarrow f(x) \leqslant f(x + \Delta x) \Rightarrow f'(x) = f'_-(x) = \lim\limits_{\Delta x \to 0-}{\dfrac{f(x + \Delta x) - f(x)}{\Delta x}} \leqslant 0$

\end{itemize}

Таким образом, $\forall x \in (a, b) \Rightarrow f'(x) \leqslant 0$.

$(\Leftarrow)$

По условию во всех точках интервала $(a,b)$, в которых $f(x)$ дифференцируема, выполняется неравенство $f'(x) \leqslant 0$. Пусть $x_1$ и $x_2,\ a < x_1 < x_2 < b,$ — произвольные точки этого промежутка. Тогда функция $y = f(x)$ удовлетворяет теореме Лагранжа, $\Rightarrow  \exists c \in (x1, x2)$, такая, что $f(x_2) - f(x_1) = f'(c)(x_2 - x_1)$

Т.к. $\forall x \in (a, b) \Rightarrow f'(x) \leqslant 0$, и $x_2 > x_1$, то $f(x_2) - f(x_1) \leqslant 0 \iff f(x_2) \leqslant  f(x_1)$

А значит, функция $y = f(x)$ - неубывающая на $(a, b)$ по определению.

Теорема доказана.
% проверить 
\newpage 
\subsection*{45. \textit{Первое достаточное условие экстремума (по первой производной)}}\addcontentsline{toc}{subsection}{45. \textit{Первое достаточное условие экстремума (по первой производной)}}
\begin{Quote2} 
\small\centering 

Используются определения 54, 64, 76, 77 из списка \end{Quote2} 

Пусть функция $y = f(x)$ непрерывна в $U_\delta(x_0)$ и дифференцируема в $\overset{\circ}U_\delta(x_0)$ Тогда, если $f'(x)$ меняет знак с минуса на плюс при переходе через точку $x_0$, то в этой точке функция $f(x)$ имеет строгий локальный минимум, а если $f'(x)$ меняет знак с плюса на минус при переходе через $x_0$, то функция $f(x)$ имеет в этой точке строгий локальный максимум. Если же $f'(x)$ сохраняет знак в проколотой окрестности точки $x_0$, то экстремума в этой точке нет.
\vspace*{20pt} 

\textit{Доказательство}

Рассмотрим первое утверждение теоремы. Если $f'(x) < 0$ при всех $x \in (x_0 - \delta, x_0)$, то на полуинтервале $(x_0 - \delta, x_0]$ функция $f(x)$ убывает, и для любого $x \in (x_0 - \delta, x_0)$ имеем $f(x) > f(x_0)$. На полуинтервале $[x_0, x_0 + \delta)$ функция $f(x)$ возрастает, и $f(x_0) < f(x)$ для всех $x \in (x_0, x_0 + \delta)$. Мы видим, что $x_0$ и в самом деле есть точка строгого локального минимума. Аналогично доказывается и второе утверждение теоремы. В случае последнего утверждения функция $f(x)$ либо возрастает, либо убывает на интервале $(x0 - \delta, x0 + \delta)$ в зависимости от знака производной $f'(x)$; экстремума в точке $x_0$ в обоих случаях нет.

Теорема доказана.
\newpage 
\subsection*{46. \textit{Второе достаточное условие экстремума (по второй производной)}}\addcontentsline{toc}{subsection}{46. \textit{Второе достаточное условие экстремума (по второй производной)}}
\begin{Quote2} 
\small\centering 

Используются определения 64, 74, 75, 81 из списка \end{Quote2} 

Если точка $x = x_0$ - стационарная точка функции $y = f(x)$, а функция $y = f(x)$ дважды дифференцируема в $x = x_0$ и $f''(x_0) > 0 \ \big(f''(x_0) < 0\big)$, тогда точка $x = x_0 $ - точка локального минимума (максимума).
\vspace*{20pt} 

\textit{Доказательство}

По условию $f''(x_0)>0 \Rightarrow$ функция $f'(x_0)$ является возрастающей в $U(x_0)$

По условию $x_0 $ - стационарная точка функции $y = f(x) \overset{\text{по опр.}}{\Rightarrow} f'(x_0)= 0$ Таким образом, $$
\begin{rcases}
    f'(x) < 0 \text{ при } x < x_0\\
    f'(x) > 0 \text{ при } x > x_0
\end{rcases}
\Rightarrow \text{по первому достаточному условию экстремума} $$

$$
x_0 - \text{ точка локального минимума}
$$

Для $f''(x_0)<0$ аналогично.

Теорема доказана.
\newpage 
\subsection*{47. \textit{Достаточное условие выпуклости функции}}\addcontentsline{toc}{subsection}{47. \textit{Достаточное условие выпуклости функции}}
\begin{Quote2} 
\small\centering 

Используются определения 64, 86, 87 из списка \end{Quote2} 

Пусть функция $f(x)$ дважды дифференцируема на интервале $(a, b)$, причем в каждой точке $x \in (a, b)$ выполняется неравенство $f''(x) > 0$. Тогда функция $f(x)$ выпукла вниз на указанном интервале. Если же во всех точках интервала $(a, b)$ вторая производная $f''(x)$ отрицательна, то функция f(x) выпукла вверх на этом интервале.
\vspace*{20pt} 

\textit{Доказательство}

Докажем лишь первое утверждение теоремы (второе доказывается аналогично). Рассмотрим касательную к графику функции $y = f(x)$ в точке $(x_0, f(x_0))$, $x_0 \in (a, b)$. Уравнение такой касательной имеет вид $y = f(x_0) + f'(x_0)(x - x_0)$. Пусть для определенности $x_0 < x < b$. Тогда разность ординат точки касательной $\big(x, f(x_0) + f'(x_0)(x - x_0)\big)$ и точки графика $\big(x, f(x)\big)$ равна $\Delta y = f(x_0) - f(x) + f'(x_0)(x - x_0).$ По теореме Лагранжа $f(x) - f(x_0) = f'(c)(x - x0).$ Поэтому $\Delta y = \big(f'(x_0) - f'(c)\big)\cdot(x - x_0), \ c \in (x_0, x).$ Применим еще раз теорему Лагранжа: $\Delta y = -f''(c_1)(c - x_0)(x - x_0),\ c_1 \in (x_0, c).$ Здесь $f''(c_1) > 0,\ c - x_0 > 0, \ x - x_0 > 0$, поэтому $\Delta y < 0$, и точка касательной лежит ниже соответствующей точки графика функции. Аналогично можно доказать это утверждение и в случае $a < x < x_0$. Таким образом, точки касательной лежат ниже соответствующих точек графика функции, и функция $f(x)$ выпукла вниз на интервале $(a, b)$. Теорема доказана.
\newpage 
\subsection*{48. \textit{Необходимое условие точки перегиба}}\addcontentsline{toc}{subsection}{48. \textit{Необходимое условие точки перегиба}}
\begin{Quote2} 
\small\centering 

Используются определения 64, 88 из списка \end{Quote2} 

Пусть функция $f(x)$ дважды дифференцируема в окрестности точки $x_0$, причем вторая производная непрерывна в указанной точке. Тогда если $x_0$ - точка перегиба графика функции $y = f(x)$, то $f''(x_0) = 0.$
\vspace*{20pt} 

\textit{Доказательство}

Предположим, $f''(x_0) \neq 0$, и пусть для определенности $f''(x_0) > 0$. Тогда в силу непрерывности $f
''(x)$ в точке $x_0$ существует окрестность $U_\delta(x_0)$ этой точки такая, что $f
''(x) > 0$ во всех точках этой окрестности по теореме о сохранении знака непрерывной функции в окрестности точки. Тогда на обоих интервалах $(x_0 - \delta, x_0)$ и $(x_0, x_0 + \delta)$ функция $f(x)$ выпукла вниз, что противоречит наличию перегиба в точке $x_0$. Поэтому предположение $f''(x_0) \neq 0$ неверное, $\Rightarrow f''(x_0) = 0$.

Теорема доказана.
% уточнить (это первое достаточное условие, есть еще второе, с производной третьего порядка)
\newpage 
\subsection*{49. \textit{Достаточное условие точки перегиба}}\addcontentsline{toc}{subsection}{49. \textit{Достаточное условие точки перегиба}}
\begin{Quote2} 
\small\centering 

Используются определения 50, 88 из списка \end{Quote2} 

Пусть функция $f(x)$ определена в окрестности $U_\delta(x_0)$ точки $x_0$ и непрерывна в указанной точке. Тогда, если в соответствующей проколотой окрестности $\overset{\circ}U_\delta(x_0)$ функция $f(x)$ имеет вторую производную, которая меняет знак при переходе через точку $x_0$, то точка $x_0$ есть точка перегиба функции $y = f(x)$.
\vspace*{20pt} 

\textit{Доказательство}

Пусть для определенности вторая производная $f
''(x)$ положительна при $x \in (x_0 - \delta, x_0)$ и отрицательна при $x \in (x_0, x_0 + \delta)$. Тогда на $(x_0 -\delta, x_0)$ функция $f(x)$ выпукла вниз, а на $(x_0, x_0 + \delta)$ выпукла вверх, т.е. при переходе через точку $x_0$ направление выпуклости меняется на противоположное. Отсюда следует, что $x_0$ - точка перегиба функции $f(x)$.

Теорема доказана.
\newpage\let\clearpage\relax 

\end{document}
    